
\documentclass[colorinlistoftodos,article,11pt,reqno]{article}
\usepackage{articlemacro}
\usepackage{tabularx}
\usepackage{etoc}


%for WIP
%\usepackage[notcite,notref]{showkeys}

\newcommand{\hldahli}[1]{{\sethlcolor{pink}\hl{#1}}\sethlcolor{yellow}}
\newcommand{\hlcan}[1]{{\sethlcolor{green}\hl{#1}}\sethlcolor{yellow}}



%========== Special macros for the article ========================


% Categories
\newcommand{\RigSH}{\textup{RigSH}}
\newcommand{\RigH}{\textup{RigH}}
\newcommand{\rig}{\textup{rig}}
\newcommand{\Sm}{\textup{Sm}}
\newcommand{\Rig}{\textup{Rig}}
\newcommand{\AffRig}{\textup{AffRig}}
\newcommand{\SmRig}{\textup{SmRig}}
\newcommand{\FSch}{\textup{FSch}}
\newcommand{\Pro}{\textup{Pro}}
\newcommand{\Cond}{\textup{Cond}}
\newcommand{\eff}{\textup{eff}}
\newcommand{\Ani}{\textup{Ani}}
\newcommand{\Spc}{\textup{Spc}}
\newcommand{\PSh}{\textup{PSh}}
\newcommand{\Sh}{\textup{Sh}}
\newcommand{\ProFin}{\textup{ProFin}}
\newcommand{\EDS}{\textup{EDS}}
\newcommand{\Cat}{\textup{Cat}}


% Superscripts and Subscripts
\newcommand{\qcqs}{\textup{qcqs}}
\newcommand{\Nis}{\textup{Nis}}
\newcommand{\an}{\textup{an}}
\newcommand{\ft}{\textup{ft}}
\newcommand{\lft}{\textup{lft}}

% K-theory
\newcommand{\K}{\textup{K}}
\newcommand{\KH}{\textup{KH}}
%\newcommand{\kk}{\textup{k}}
\newcommand{\Kan}{\K^\textup{an}}
\newcommand{\kan}{\textup{k}^\textup{an}}
\newcommand{\Kcont}{\K^\textup{cont}}
\newcommand{\kcont}{\textup{k}^\textup{cont}}
\newcommand{\KHcont}{\textup{KH}^\textup{cont}}
\newcommand{\khcont}{\textup{kh}^\textup{cont}}
\newcommand{\KVan}{\textup{KV}^\textup{an}}

% Other operators
\newcommand{\Map}{\textup{Map}}
\newcommand{\Mapu}{\underline{\Map}}
 
% Rigid Geometry
\newcommand{\Arig}{\AA^{1}}
\newcommand{\Anrig}{\AA^{n}}
\newcommand{\Brig}{\BB^{1}}
\newcommand{\Bnrig}{\BB^{n}}
\newcommand{\skp}[1]{\langle #1 \rangle}

% Arrows and decorations
\renewcommand{\to}[1][]{\overset{#1}{\rightarrow}}	
\newcommand{\To}[1][]{\overset{#1}{\longrightarrow}}	
\newcommand{\inj}[1][]{\overset{#1}{\hookrightarrow}}		% injective maps
\newcommand{\surj}[1][]{\overset{#1}{\twoheadrightarrow}}		% surjective maps


\DeclareMathOperator*{\prolim}{``lim''}
\renewcommand{\L}{\textup{L}}
\newcommand{\Lmot}{\L_{\textup{mot}}}
\newcommand{\es}{\emptyset}
\newcommand{\PrL}{\textup{Pr}^{\textup{L}}}
\newcommand{\PrLO}{\textup{Pr}^{\textup{L},\otimes}}
\newcommand{\CAlgPr}{\textup{CAlg}(\PrLO)}

% Abbrevations
\newcommand{\infcat}{category}
\newcommand{\infcats}{categories}
\newcommand{\inftop}{$\infty$-topos}
\newcommand{\inftops}{$\infty$-topoi}

%\newcommand{\Xcal}{\mathcal{X}}
%\newcommand{\Ycal}{\mathcal{Y}}
%\newcommand{\Zcal}{\mathcal{Z}}

% Dahlis Diagramme (irgendwann sollte ich voll auf tikzd umsteigen...)
\newcommand{\carre}[4]{\[\begin{xy}\xymatrix{#1\ar[r]\ar[d]&#2\ar[d]\\#3\ar[r]&#4}\end{xy}\]}
\newcommand{\carremap}[8]{\[\begin{xy}\xymatrix{#1\ar[r]^{#2}\ar[d]_{#3}&#4\ar[d]^{#5}\\#6\ar[r]_{#7}&#8}\end{xy}\]}
\newcommand{\carremaptag}[9]{\begin{align}\begin{xy}\tag{#1}\xymatrix{#2\ar[r]^{#3}\ar[d]_{#4}&#5\ar[d]^{#6}\\#7\ar[r]_{#8}&#9}\end{xy}\end{align}}

\newcommand\blfootnote[1]{%
  \begingroup
  \renewcommand\thefootnote{}\footnote{#1}%
  \addtocounter{footnote}{-1}%
  \endgroup
}


%==================================================================

\date{}
\title{Towards $\mathbb{A}^1$-homotopy theory of rigid analytic spaces}
\author{Christian Dahlhausen and Can Yaylali}


%\usepackage{authblk}
%\author[1]{Christian Dahlhausen}
%\author[2]{Can Yaylali}
%\affil[1]{Institut für Mathematik, Im Neuenheimer Feld 205, 69120 Heidelberg, Germany}
%\affil[2]{Institut de Mathématique d'Orsay, 307 Rue Michel Magat Bâtiment 307, 91400 Orsay, France}
%\email{}

\begin{document}
\maketitle
\begin{abstract}
%Foonote with MSC classifications
\let\thefootnote\relax\footnote{2020 Mathematics Subject Classification: 14F42, 19E99, 14G22}

To any rigid analytic space (in the sense of Fujiwara-Kato) we assign an $\mathbb{A}^1$-invariant rigid analytic homotopy category with coefficients in any presentable category. We show some functorial properties of this assignment as a functor on the category of rigid analytic spaces. Moreover, we show that there exists a full six functor formalism for the precomposition with the analytification functor by evoking Ayoub's thesis. As an application, we identify connective analytic K-theory in the unstable homotopy category with both $\mathbb{Z}\times\textup{BGL}$ and the analytification of connective algebraic K-theory. As a consequence, we get a representability statement for coefficients in light condensed spectra.
\end{abstract}



\thispagestyle{empty}
\setcounter{tocdepth}{1}
\tableofcontents
\setcounter{tocdepth}{2}


\section{Introduction}
The $\AA^1$-homotopy category of schemes was constructed and studied by Morel\hyp{}Voevodsky \cite{MV1}.
In modern language, their homotopy category $\textup{H}(S)$ is the category of $\AA^1$\hyp{}invariant Nisnevich sheaves on smooth schemes over a base scheme $S$ with values in spaces. Under this construction they show, under some regularity assumptions, that connective algebraic K-theory can be represented as $\ZZ\times \BGL\in \textup{H}(S)$.
Moreover, after $\otimes$-inverting $\PP^1$, we obtain the stable homotopy category $\textup{SH}(S)$. In his thesis, Ayoub showed that the assignment $S\mapsto \textup{SH}(S)$ admits a full six functor formalism \cite{AyoubThesis}. 
Ayoub deduces this from properties of the functor $\textup{SH}(-)$ such as $\AA^1$-invariance, $\PP^1$-stability, and the existence of the localisation sequence. In this article, we are going to extend these results to the analytic setting, using the rigid analytic line $\Arig$ as an interval.


\vspace{6pt}\noindent Motives on rigid analytic varieties have been constructed and studied by Ayoub\linebreak \cite{AyoubRig}. The idea is to follow Morel-Voevodsky's approach in constructing the stable homotopy category $\textup{RigSH}^{\Brig}(X)$ for a rigid analytic variety $X$. Ayoub defines the category of rigid motives by Nisnevich localisation, contracting the closed unit disc $\Brig$, and then $\otimes$-inverting the unit disc without the origin. In \textit{op.\!\! cit.}\ it is shown that this definition yields a coefficient system in the sense that any morphism $f$ of rigid analytic varieties induces a colimit preserving functor $f^*$ on $\textup{RigSH}^{\Brig}(-)$ with the following properties.
\begin{enumerate}
    \item[(PF)] If $f$ is smooth, then $f$ admits a left adjoint satisfying base change and projection formula.
    \item[(Loc)] Any closed immersion $i\colon Z\hookrightarrow X$ of rigid analytic varieties with open complement $j \colon U\hookrightarrow X$ induces a fibre sequence $j_\#j^* \to \id \to i_*i^*$ of endofunctors on $\textup{RigSH}^{\Brig}(X)$.
\end{enumerate}
 While the proof of (PF) is a formality, the proof of (Loc) needs more care and follows the idea of Voevodsky. These two properties and the definition immediately imply that on algebraic maps $X\rightarrow Y$ of rigid analytic varieties there exists a six functor formalism \cite[Scholie 1.4.2]{AyoubThesis}. More recently, in joint work Ayoub-Gallauer-Vezzani generalise this result to arbitrary finite type maps of rigid analytic varieties \cite{AGV}.
\par
A drawback of this definition of $\textup{RigSH}^{\Brig}$ is that not all cohomology theories on rigid analytic varieties satisfy $\Brig$-homotopy invariance. For instance, continuous K-theory $\Kcont$ defined by Morrow \cite{Morrow}  and analytic K-theory defined by Kerz-Saito-Tamme \cite{kst-i} are not $\Brig$-invariant (Remark~\ref{Kan-not-B1-invariant--remark}). Hence, these theories cannot be represented in $\textup{RigDA}$. On the other hand, letting $\Arig$ be the rigid affine line, analytic K-theory is $\Arig$-invariant \cite[Cor.~2.7]{kst-ii} and  -- assuming resolution of singularities -- continuous K-theory is $\Arig$-invariant for regular rigid analytic varieties \cite[Prop. 5.14]{kst-i}.
Doing motivic homotopy theory with $\Arig$ as the interval has already been done by Sigloch \cite{sigloch-thesis}.
This suggests that one should define the category of motives on a rigid analytic variety $X$ as 
\[
    \RigSH(X)\coloneqq \Sh_{\Nis}^{\AA^1}(\SmRig_{X})[\textup{Th}(\Brig)^{-1}],
\]
where $\textup{Th}(\Brig)$ denotes the Thom motive of the rigid closed unit disc. This is equivalent to $\otimes$-inverting $\PP^1$, see section~\ref{sec.thom-motives-stabilisation}.


\vspace{6pt}\noindent 
Eventually, we apply our theory to K-theory of rigid analytic spaces, namely analytic K-theory as defined and studied by Kerz-Saito-Tamme \cite{kst-i,kst-ii}. Since this K-theory takes values in pro-categories, it makes sense to consider motives with coefficients in an arbitrary symmetric monoidal stable presentable \infcat{} $\Vcal_{\textup{st}}$. Then we can apply the theory to the categories of pro-spectra and, since pro-categories do not behave that nicely, (light or $\kappa$-small) condensed spectra.

Let us fix a symmetric monoidal presentable category $\Vcal$ and its stabilisation $\Vcal_{\textup{st}}$. The motivating examples include the category $\Cond_{\kappa}(\Spc)$ of $\kappa$-small condensed spaces and the category $\Cond_{\kappa}(\Sp)$ of $\kappa$-small condensed spectra, for a fixed uncountable strong limit cardinal cardinal $\kappa$. The category $\RigSH(-,\Vcal)$ keeps the property (PF) above and perhaps more surprisingly also (Loc), even without stabilising. 

\begin{intro-theorem}[\protect{\ref{prop-smooth-BC-and-PF},~\ref{thm-gluing}}]
The assignment $X\mapsto \Sh^{\Arig}_{\Nis}(X,\Vcal),\ f\mapsto f^{*}$, from rigid analytic varieties to symmetric monoidal presentable $\Vcal$-linear categories satisfies (PF) and (Loc). 
\end{intro-theorem}

As an immediate consequence, we get a partial six functor formalism on $\RigSH$ over a nonarchimedean field $k$.

\begin{intro-corollary}[\protect{\ref{thm-6ff}}]
   The assignment $X\mapsto \RigSH(X^{\an},\Vcal_{\textup{st}}),\ f\mapsto f^{*}$, from separated finite type $k$-schemes to symmetric monoidal stable presentable $\Vcal_{\textup{st}}$-linear categories admits a six functor formalism satisfying base change, purity, and projection formula.
\end{intro-corollary}

An extension of the six functor formalism to non-algebraic maps of rigid analytic varieties is not clear. In the classical rigid motivic theory of Ayoub, the rigid affine line and the closed unit disc are $\Brig$-homotopic \cite[Prop. 1.3.4]{AyoubRig}. This will not be the case, when working with $\Arig$-homotopy instead. This shows that we need more careful treatment to obtain a full six functor formalism for rigid spaces and cannot rely purely on the results of \cite{AGV}. Nevertheless, in future work we aim to establish a such full six functor formalism.

\vspace{6pt}\noindent 
Let us pass to an application of our theory. By design, analytic K\hyp{}theory of rigid spaces is an $\Arig$\hyp{}invariant Nisnevich sheaf, hence it yields an object of the unstable $\Arig$\hyp{}homotopy category. Assuming resolution of singularities, we get the following identification (which we state in the introduction in less generality than in the main text).

\begin{intro-theorem}[{\ref{thm:representability}}] 
\label{thm:representability-intro}
Let $k$ be a discretely valued nonarchimedean field and assume $(\dagger)_k$, see section~\ref{assumption-dagger}. Then there is a canonical equivalence
\[ \Omega^\infty\Kan_{\geq 0} \simeq \ZZ\times\BGL \]
in the category $\RigH(k)$.
In particular, for every rigid space $X$ there is a functorial equivalence
\[ \Omega^\infty\Kan_{\geq 0}(X) \simeq \Homline_{\Cond^\omega(\Spc)}(\Lmot X,\Lmot(\ZZ\times\BGL)) \]
in the category $\Cond^\omega(\Spc)$; here the right-hand side denotes the enriched Hom-space as a left $\Cond^\omega(\Spc)$-module, see Appendix~\ref{app.yoneda}.
\end{intro-theorem}

This is analogous to the identification $\K_{\geq 0} \simeq \ZZ\times\BGL$ of algebraic K-theory in the Morel-Voevodsky category $\textup{H}(k)$ \cite[Prop.~3.10]{MV1}. In particular, assuming $(\dagger)_k$, the analytification functor $\textup{H}(k)\to\RigH(k)$ maps connective algebraic K-theory to connective analytic K-theory. See also \cite[p.~ix]{AyoubRig}.

We expect Theorem~\ref{thm:representability-intro} to be helpful in the study of analytic K-theory whose underlying spectrum is related to  Efimov's K-theory of dualisable categories and Clausen-Scholze's nuclear modules on analytic spaces, see Remark~\ref{connection-efimov-k-theory--rem}.

Furthermore, we prove an $\Arig$-analytic Bass Fundamental Theorem for analytic K-theory. Its proof relies on the analogous result for $\Brig$ by Kerz-Saito-Tamme \cite[Cor.~2.6]{kst-ii}.
\begin{intro-theorem}[{\ref{bass-fundamental-theorem-for-Kan--cor}}]
For every $n\in\ZZ$ and every $X\in\Rig_k$ there is an exact sequence
\[ 0 \To \Kan_n(X) \To \Kan_n(X\times\Arig) \oplus \Kan_n(X\times\AA^{-1}) \To[\pm] \Kan(X\times\Gm) \To[\partial] \Kan_{n-1}(X) \to 0 \]
of pro-abelian groups where the map $\partial$ has a split.
\end{intro-theorem}

Using the Bass Fundamental Theorem together with Theorem~\ref{thm:representability-intro} we will see that the analytification functor maps the $\PP^1$-spectrum  representing algebraic K-theory to a $\PP^1$-spectrum representing analytic K-theory under assumption $(\dagger)_k$. 

\begin{intro-corollary}[\protect{\ref{cor-K-P1-spectrum}}]
There exists a $\PP^1$-action on $\ZZ\times\BGL\in\RigH(k)$ inducing a $\PP^1$-spectrum $\textup{KGL}^{\an}\in \RigSH(k)$. Under assumption $(\dagger)_k$, the $\PP^1$\hyp{}spectrum $\textup{KGL}^{\an}$ represents (non\hyp{}connective) analytic $K$\hyp{}theory.
\end{intro-corollary}


\vspace{6pt}\noindent
Weibel defined homotopy K-theory as the $\Arig$-localisation of algebraic K-theory. In this spirit, we define \emph{continuous homotopy K-theory} $\KHcont$ (Defintion~\ref{defi-continuous-homotopy-k-theory}) and identify it with analytic K-theory (Theorem~\ref{A1-localisation--thm}) by evoking a result of Kerz-Saito-Tamme.



\vspace{6pt}\noindent\textbf{Acknowledgements.}
We heartily thank Joseph Ayoub for guiding us through his work on six functor formalisms \cite{AyoubThesis} and rigid analytic motives \cite{AyoubRig} and Martin Gallauer for insightful explanations about his work \cite{AGV} and Brad Drew's work on coefficient systems \cite{DrewMHM}. We thank Peter Scholze for communicating to us the proof of Theorem~\ref{comparison-conservative--thm} and Vlaidimir Hinich for answering our questions about the enriched Yoneda lemma.
We thank Yicheng Zhou (Lemma~\ref{pro-vanishing-and-Mittag-Leffler--lemma}) and Marc Hoyois (Theorem~\ref{thm-gluing}) for helpful comments on a previous version of this article.
Furthermore, we thank Ryomei Iwasa and Florian Strunk for helpful discussions around this paper's content. The first named author thanks Georg Tamme for continuous discussions about K-theory of rigid spaces.

Both authors were supported by the Deutsche Forschungsgemeinschaft (DFG)\linebreak through the Collaborative Research Centre TRR 326 \textit{Geometry and Arithmetic of Uniformized Structures}, project number 444845124. The second named author was furthermore supported through the Walter Benjamin Programme (DFG, project number 524431573) during his stay at the Laboratoire de Mathématiques d'Orsay. He wants to thank Vincent Pilloni for hosting him and all the other people giving a friendly and mathematically interesting environment.

\vspace{6pt}\noindent\textbf{Assumptions and notations.}
Throughout, we fix some inaccessible regular cardinal $\kgbar$. By \textit{small}, we will mean \textit{$\kgbar$-small}. 
A \emph{category} always means an $\infty$-category, i.e. an $(\infty,1)$-category in the sense of \cite{HTT}, and we identify 1-categories with their image under the nerve functor.
%We will freely work with $(\infty,1)$-categories in the sense of \cite{HTT} and if we write category, we will always mean $(\infty,1)$-category. When we say \textit{category}, we always mean $1$-category and view it as a category via the nerve functor.\par
Let us fix some notation for the categories appearing in this article.
\begin{enumerate}%[\itemsep=2em]
    \item[$\bullet$] $\Spc$, the category of small spaces, or equivalently small $\infty$-groupoids or small $\infty$-sets, or small anima (depending on your taste).
    \item[$\bullet$] $\Sp$, the category of spectra.
    \item[$\bullet$] $\PrL$, the category of presentable categories with colimit preserving functors.
    \item[$\bullet$] $\PrLO$, the category $\PrL$ endowed with the Lurie tensor product \cite[\S 4.8.1]{HA}.
\end{enumerate}
The term \emph{essentially unique} is short hand for the expression \emph{unique up to contractible choice}.






%%%%%%%%%%%%%%%%%%%%%%%%%%%%%%%%%%%%%%%%%%%%%%%%%%%%%%
\section{Background in rigid analytic geometry}
%%%%%%%%%%%%%%%%%%%%%%%%%%%%%%%%%%%%%%%%%%%%%%%%%%%%%%

We place ourselves in the same setting as Ayoub-Gallauer-Vezzani \cite[\S 1]{AGV} and use the notion of \emph{rigid spaces}\footnote{For the sake of brevity, we mostly use the terminology `rigid space' instead of the more accurate `rigid analytic space'.} as defined by Fujiwara-Kato \cite{fuji-kato} which goes back to Raynaud \cite{raynaud} and was also developed by Abbes \cite{egr}. 

\vspace{6pt}\noindent\textbf{Rigid spaces.} Let $\FSch$ be the category of formal schemes and let $\FSch^\qcqs$ be its full subcategory of quasi-compact and quasi-separated formal schemes. The category\linebreak $\Rig^\qcqs$ of quasi-compact and quasi-separated \emph{rigid spaces} is defined to be the 1\hyp{}categorical localisation of $\FSch^\qcqs$ by the set of admissible formal blow-ups. The  category $\Rig$ is  defined by gluing quasi-compact and quasi-separated rigid spaces along open immersions, see \cite[ch.~II, \S 2.2~(c)]{fuji-kato}. For a formal scheme $\Xcal$ denote by $\Xcal^\rig$ its associated rigid space.

\vspace{6pt}\noindent\textbf{Relation to adic spaces.} There exists a fully faithful functor $\textup{Adic}^\textup{unif}\inj\Rig$ from the category of uniform adic spaces into the category of rigid spaces sending $\Spa(R,R^+)$ to $\Spf(R^+)^\rig$ which is compatible with gluing along open immersions \cite[\S 1.2]{AGV}.
We call a rigid space \emph{adic} if it lies in the essential image of this functor.

\vspace{6pt}\noindent\textbf{Morphisms locally of finite type.}
A morphism $f\colon Y\to X$ of rigid spaces is said to be \emph{locally of finite type} if there exists an open cover $(X_i)_i$ of $X$ such that the restricted morphisms $f^{-1}(X_i)\to X$ admit formal models which are locally of finite type, i.e.\ locally given by algebras of topologically finite type, see  \cite[ch.~II, sec.~2.3]{fuji-kato}.
Note that an algebra of topologically finite type over a uniform Huber ring, is again a uniform Huber ring. It follows that, if $X$ is adic, then $Y$ is adic as well.

\vspace{6pt}\noindent\textbf{Smooth morphisms.}
A morphism $f\colon Y\to X$ of rigid spaces is called \emph{smooth} if, locally on $Y$ and $X$, there exists a formal formal model $\Ycal\to\Xcal$ which, locally on rings, is given as a composition $A\to A\langle t_1,\ldots,t_n\rangle \to[\phi] B$ where $\phi$ is rig-\'{e}tale, see \cite[1.3.3., 1.3.13]{AGV}.
It holds true that a morphism of adic rigid spaces is smooth in this sense if and only if it smooth in the sense of Huber, see \cite[1.3.14]{AGV} and the references there.

\vspace{6pt}\noindent\textbf{The Nisnevich topology.}
A family $(\Ycal_i\to\Xcal)_i$ in $\FSch$ is a \emph{Nisnevich cover} if the induced family $(\Ycal_{i,\sigma}\to \Xcal)_i$ on special fibres is a Nisnevich cover of schemes. A family $(Y_i\to X)_i$ in $\Rig$ is a \emph{Nisnevich cover} if locally on $X$ the family admits a refinement that has a formal model which is a Nisnevich cover of formal schemes.

\vspace{6pt}\noindent\textbf{The rigid affine line.}
Let us quickly recall some background concerning the rigid affine line, for more details consider the lecture notes of Scholze-Weinstein \cite[Lecture~4]{weinstein-scholze} or Hübner \cite[Ex.~1.4.3, Ex.~1.10.4]{huebner-adic}. 
Let $(R,R^+)$ be a uniform Huber pair and set $S=\Spa(R,R^+)$.
Throughout this article, we work with the rigid affine line, seen as a rigid space\footnote{We will mostly work in the setting of rigid spaces and therefore omit a superscript indicating that $\Arig$ is not the schematic affine line.}
\[
    \Arig_{S} \coloneqq \Spa(\ZZ[T],\ZZ) \times S
\]
using the fully faithful embedding $\textup{Adic}^\textup{unif} \inj \Rig$ from above.\footnote{The ring $\ZZ[T]$ is equipped with the discrete topology. For the existence of fibre products we refer to \cite[ch.~II, sec.~2.4]{fuji-kato}.} This description of the affine line will make it easy to construct explicit homotopies. All of the computations (e.g.\ descriptions of homotopies) can be reduced to computations on $\Spa(\ZZ[T],\ZZ)$. 
\par
Analogously, we define the closed rigid unit disc as 
\[
    \Brig_{S} \coloneqq \Spa(\ZZ[T],\ZZ[T]) \times S = \Spa(R\langle T\rangle,R^+\langle T\rangle)
\]
where $R\langle T\rangle$ denotes the ring of those power series in $T$ whose coefficients tend to $0$. 

Now assume that $R$ is a Tate ring and let $\pi$ be a pseudo-uniformiser. Then we can describe the affine line as the union of closed units discs with larger and larger radius
\begin{equation}
    \label{eq-Arig-colim}
    \Arig_{S} = \bigcup_{n\geq 0} \Spa(R\langle \pi^n T\rangle,R^+\langle \pi^n T\rangle).
\end{equation}

    


%%%%%%%%%%%%%%%%%%%%%%%%%%%%%%%%%%%%%%%%%%%%%%%%%%%%%%
\section{Unstable  \texorpdfstring{$\Arig$}~-homotopy theory in rigid geometry}

\label{sec.unstable}
\localtableofcontents
%%%%%%%%%%%%%%%%%%%%%%%%%%%%%%%%%%%%%%%%%%%%%%%%%%%%%%
\vspace{6pt}
\noindent In this section, we want to analyse $\Arig$-invariant Nisnevich sheaves on smooth rigid spaces, also called unstable rigid homotopy category. The upshot is that this category will define a pullback formalism in the sense that for any smooth morphism of rigid spaces $f$ the pullback $f^*$ admits a left adjoint satisfying base change and projection formula. But we will go even a bit further and prove that we even have a localisation sequence in this more general setting.

\vspace{6pt}\noindent\textbf{Notation.} Throughout this section, let $(R,R^+)$ be a uniform Huber pair with associated rigid space $S=\Spa(R,R^+)$ and let $\Dcal$ be a presentable category of coefficients.

%%%%%%%%%%%%%%%%%%%%%%%%%%%%%%%%%%%%%%%%%%
\subsection{Homotopy invariant presheaves on rigid spaces}
%%%%%%%%%%%%%%%%%%%%%%%%%%%%%%%%%%%%%%%%%%

As mentioned earlier, there are several reasonable choices of an interval object to do homotopy theory with rigid spaces such as the closed unit disc $\Brig$ and the rigid affine line $\Arig$. If the target category is a pro-category, one can also speak about pro-$\Brig$-invariance (Lemma~\ref{pro-invariant=Aan-invariant--lem}).

\begin{defi} \label{Aan-B1-invariance--defin}
We say that a presheaf $F\in\Fun(\Rig_S\op,\Dcal)$ is \ldots
\begin{enumerate}
	\item $\Brig$\emph{-invariant} if for every $X\in\Rig_S$ the canonical map $F(X) \to F(X\times\Brig)$ is an equivalence in $\Dcal$.
	\item $\Arig$\emph{-invariant} if for every $X\in\Rig_S$ the canonical map $F(X)\to F(X\times\Arig)$ is an equivalence in $\Dcal$.
\end{enumerate}
\end{defi}

\begin{defi}
Let $\Ccal$ be a site. A presheaf $F\in\Fun(\Ccal\op,\Dcal)$ is said to be a \emph{sheaf} if for every object $X$ of $\Ccal$ and every covering sieve $U\inj X$ the canonical map $F(X)\to F(U)$ is an equivalence in $\Dcal$; here $F(U)$ is defined via the equivalence
	\[
	\Fun^\mathrm{lim}(\PSh(\Ccal)\op,\Dcal) \To \Fun(\Ccal\op,\Dcal).
	\]
Thus if we write $U \simeq \colim_iU_i$ as a colimit of representables, then $F(U) \simeq \lim_iF(U_i)$.
\end{defi}

\begin{lemma} \label{B1-implies-Aan-invariance--lem}
Every $\Brig$-invariant sheaf on $\Rig_S$ is $\Arig$-invariant.
\end{lemma}
\begin{proof}
We can identify $\Arig$ with the colimit $\colim_{t\mapsto\pi t}\Brig$ in the category $\Rig_S$ where $t$ is a parameter for $\Brig=\Spa(k\skp{t}),k^\circ\skp{t})$. Since $F$ is a sheaf (and hence preserves this colimit) and since a colimit of open immersions commutes with fibre products, we have for every sheaf $F$ equivalences
	\begin{align*}
	F(X\times\Arig) \simeq F(X\times\colim_{t\mapsto\pi t}\Brig)
	\simeq F(\colim_{t\mapsto\pi t}(X\times\Brig))
	\simeq \lim_{t\mapsto\pi t} F(X\times\Brig).
	\end{align*}
Hence the map $F(X) \to F(X\times\Arig)$ is an equivalence if the map $F(X) \to F(X\times\Brig)$ is an equivalence.
\end{proof}


\begin{defi}
Denote by $\PSh^{\Arig}(\Rig_k,\Dcal)$ the full subcategory of the category of\linebreak presheaves $\PSh(\Rig_k,\Dcal)$ spanned by $\Arig$-invariant presheaves.
\end{defi}

\begin{remark} \label{Kan-not-B1-invariant--remark}
There are $\Arig$-invariant sheaves which are not $\Brig$-invariant. For instance, analytic K-theory $\Kan \colon \Rig_k\op \to \Pro(\Sp^+)$ is a sheaf \cite[6.15]{kst-i} \cite[4.3]{kst-ii} which is $\Arig$-invariant \cite[4.4]{kst-ii}. However, assuming resolutions of singularities, it is not $\Brig$-invariant as it agrees with continuous K-theory \cite[6.19]{kst-i} which on $\pi_0$ is equivalent to $\K_0$ \cite[5.10]{kst-i} which has the non-$\Brig$-invariant Picard group as a direct summand \cite[Example~2]{kst-bass-quillen}. For more details about continuous K-theory and analytic K-theory we refer to subsection~\ref{subsec.continuous-analytic-k-theory}.
\end{remark}

\begin{lemma} \label{homotopy-invariantification--lem}
For every cocomplete category $\Dcal$ there exists a localisation functor
	\[
	\L_{\Arig} \colon \PSh(\Rig_S,\Dcal) \To \PSh^{\Arig}(\Rig_S,\Dcal)
	\]
which is left-adjoint to the inclusion. It has the following explicit description: for every $F\in\PSh(\Rig_k,\Ccal)$ and every $X\in\Rig_S$ there is a functorial equivalence
	\[
	(\L_\Arig F)(X) \simeq \colim_{\Delta\op} F(X\times \Delta^{\mathrm{an},\bullet})
	\]
where $\Delta^{\mathrm{an},\bullet}$ is the analytification of its algebraic analogue.
\end{lemma}
\begin{proof}
For the proof, denote by $\textup{H} F$ the presheaf whose values are the right-hand side of the desired identification. This yields a natural transformation $\alpha \colon \id \to \textup{H}$  of presheaves with values in $\Dcal$. Assume that the following conditions hold:
\begin{enumerate}
	\item[(i)] For every presheaf $F$, the presheaf $\textup{H} F$ is $\Arig$-invariant.
	\item[(ii)] For every $\Arig$-invariant presheaf $F$, the map $\alpha_F \colon F \to \textup{H} F$ is an equivalence.
\end{enumerate} 
These conditions imply for every presheaf $F$ that both maps $\textup{H}(\alpha_F),\alpha_{\textup{H} F} \colon \textup{H} F \to \textup{H}\textup{H} F$ are equivalences. Then the desired claim follows from \cite[5.2.7.4]{HTT}. Hence it suffices to check conditions (i) and (ii).

(i) Let $X\in\Rig_k$ and let $p\colon X\times\Arig \to X$ the projection and $\sigma \colon X \to X\times\Arig$ the zero section. As $p\circ\sigma=\id$ it remains to show that the induced map $p^*\circ\sigma^*$ on $(\textup{H} F)(X\times\Arig)$ is equivalent to the identity. For this purpose, we claim that the maps
	\[
	p^*\circ\sigma^*,\id \colon X\times\Arig\times\Delta^{\mathrm{an},\bullet} \rightrightarrows X\times\Arig\times\Delta^{\mathrm{an},\bullet}
	\]
are simplicial homotopic (see Remark~\ref{simplicial-homotopic--remark} below). In degree $n$, a homotopy $h_n$ on $X\times\Arig\times\Delta^{\mathrm{an},n}$ is given by $h_n(f)(t) = (\sum_{j\in f^{-1}(1)}t_j)\cdot t$ for a map $f \colon [n] \to{} [1]$. We have $i_0^*h(t)=0$ and $i_1^*h(t)=t$.

(ii) We follow the classical reasoning \cite[ch.~IV, Lemma~11.5.1]{k-book}. Let $F$ be an $\Arig$-invariant presheaf. Then for every $X\in\Rig$ there is a morphism $X\times\Delta^{\an,\bullet} \to X$ of simplicial rigid spaces which is induces by the projections. For every $n\in\Delta\op$ the induced map $F(X) \to F(X\times\Delta^{\an,n})$ is an equivalence. Hence the map
\[ F(X) = \colim_{\Delta\op} F(X) \To \colim_{\Delta\op} F(X\times\Delta^{\an,n}) = \textup{H}F(X) \]
is an equivalence. 


\begin{comment}
Indeed, for $\Dcal=\Sp$ we can follow the proof in the algebraic setting \cite[ch.ÏV, Thm.~12.3]{k-book}, there exists a spectral sequence
\[ E^1_{p,q} = \textup{N}^pF_q(X) \Rightarrow \textup{H}F_{p+q}(X) \]
where $\textup{N}^pF_q(X)$ is the intersections of the maps
\[ F_q(X\times\AA^p) \To F_q(X\times\AA^{p-1})\]
mapping $t_i\to 0$ and hitting all the other coordinates.
\end{comment}

\end{proof}

\begin{remark} \label{simplicial-homotopic--remark}
Let $\Ccal$ be an \infcat{} and $\mathrm{s}\Ccal$ the simplicial objects in $\Ccal$. We have a diagram
	\[
	\Delta \overset{i_0}{\underset{i_1}{\rightrightarrows}} \Delta_{/[1]} \to[\delta] \Delta
	\]
where $i_k([n]) = ([n]\to{}[0]\to[k][1])$ and $\delta([n]\to{}[1]) = [n]$. A \emph{simplicial homotopy} between two maps $f,g \colon X_\bullet \to Y_\bullet$ in $\mathrm{s}\Ccal$ is a map $h \colon \delta^*(X_\bullet) \to \delta^*(Y_\bullet)$ in $\Fun((\Delta_{/[1]})\op,\Ccal)$ such that $i_0^*(h)=f$ and $i_1^*(h)=g$. In this case, if $\Ccal$ is cocomplete, the induces maps $|f|,|g| \colon |X_\bullet| \to |Y_\bullet|$ on geometric realisations are homotopic in $\Ccal$. This follows since the map $\colim_{(\Delta_{/[1]})\op} \delta^*(X_\bullet) \to |X_\bullet| = \colim_{\Delta\op} X_\bullet$ is an equivalence. 
\end{remark}






%%%%%%%%%%%%%%%%%%%%%%%%%%%%%%%%%%%%%%%%%%%%%%%%%%
\subsection{The unstable analytic motivic category}
%%%%%%%%%%%%%%%%%%%%%%%%%%%%%%%%%%%%%%%%%%%%%%%%%%

\begin{defi} \label{RigH--definition}
Let $S$ be a rigid space. The \emph{unstable rigid motivic homotopy category with coefficients in $\Dcal$} is the reflective subcategory 
\[ \RigH(S,\Dcal) := \Sh^{\Arig}_\Nis(\SmRig_S,\Dcal) \]
of $\Sh_\Nis(\SmRig_S,\Dcal)$ which is spanned by $\Arig$-invariant sheaves. We have the following diagram of localisation functors:

$$ 
\begin{tikzcd}
    & \Sh_\Nis(\SmRig_S,\Dcal) \arrow[dr,"\L_{\Arig}^\#"]& \\
     \PSh(\SmRig_S,\Dcal) \arrow[ur,"\L_\Nis"] \arrow[dr,"\L_{\Arig}"] \arrow[rr,"\Lmot"] && \RigH(S,\Dcal) \\
     & \PSh^{\Arig}(\SmRig_S,\Dcal) \arrow[ur,"\L_\Nis^\#"]
\end{tikzcd} 
$$

We note that the \textit{motivic localisation} $\Lmot$ is given by the colimit of functors 
\[ \Lmot \simeq \colim( \L_\Nis\L_{\Arig} \to \L_\Nis\L_{\Arig}\L_\Nis\L_{\Arig} \to \ldots ). \]
\end{defi}





%%%%%%%%%%%%%%%%%%%%%%%%%%%%
\subsection{Analytification for motivic spaces}
%%%%%%%%%%%%%%%%%%%%%%%%%%%%


Let $k$ be a nonarchimedean field and let $S$ be a $k$-scheme of finite type. In this section, we construct an analytification functor from the algebraic stable homotopy category $\textup{H}(S,\Dcal)$ to the rigid analytic $\Arig$-homotopy category $\RigH(S^\an,\Ecal)$ depending on a \emph{change of coefficients functor} $\Dcal\to\Ecal$ between symmetric monoidal presentable \infcats. Later on, we apply this with the functors $\Spc\inj\Cond(\Spc)$ and $\Sp\inj\Cond(\Sp)$.

\vspace{6pt}\noindent\textbf{Notation.}
In this section, let $R_0$ be an adic ring with ideal of definition $I$. We set $S:=\Spec(R_0)\setminus\textup{V}(I)$ and $S^\an:=\Spf(R_0)^\rig$.
In the special case where $R_0$ is a ring of definition of a Huber pair $(R,R^+)$, we  have $S=\Spec(R)$ and $S^\an=\Spa(R)=\Spa(R,R^+)$. For any $R$-algebra $A$ we set $A^+$ as the integral closure of (the image of) $R^+$ in $A$ and write $\Spa(A):=\Spa(A,A^+)$ for simplicity.

\begin{lemma} \label{analytification-functor--lem}
There exists an analytification functor $(-)^\an \colon \Sch_S^\lft \to \Rig_{S^\an}^\lft$ sending the algebraic affine line $\Arig_S$ to the rigid analytic affine line $\Arig_{S^\an}$. In particular, for every $X\in\Sch^\lft_S$ we get an analytification functor $(-)^\an \colon \Sch_X^\lft \to \Rig_{X^\an}^\lft$. By restriction, we get a functor $(-)^\an\colon\Sm_X\to\SmRig_{X^\an}$.
\end{lemma}
\begin{proof}
For general rigid spaces the analytification functor has been constructed by\linebreak Fujiwara-Kato \cite[ch.~2, sec.~9.1.]{fuji-kato}. We give a sketch of the construction in the adic setting; for a full proof we refer to Hübner's lecture notes \cite[\S 1.10]{huebner-adic}.  

So, assuming that $S=\Spec(R)$ for a Huber pair $(R,R^+)$, the analytification functor is locally given by the assignment
\[
	R[t_1,\ldots,t_n]/(f_1,\ldots,f_k) \mapsto \bigcup_{r\geq 1} \Spa\bigl( R\skp{t_1,\ldots,t_n}_r/(f_1,\ldots,f_k)\bigr)
\]
where $R\skp{t_1,\ldots,t_n}_r$ denote the subring of $R\llbracket t_1,\ldots,t_n\rrbracket$ consisting of power series\linebreak $\sum_{i\in\NN^n}a_it^i$ such that $a_ir^i$ converges to zero as $|i|\to\infty$. This assignment is compatible with gluing and extends to the desired functor.

Alternatively, one can associate with any $X\in\Sch_S^\ft$ a natural presheaf $h_\an(X)$ on $\Rig_k$ which turns out to be representable by a rigid space $X^\an$ yielding the analytification functor, see  \cite[\S 1.1.3]{AyoubRig}.


\end{proof}

\begin{lemma} \label{analytification-presheaves--lem}
The analytification functor $(-)^\an \colon \Sch_S^\lft \To \Rig_{S^\an}^\lft$ extends to an essentially unique colimit-preserving functor
\[
\PSh(\Sch_S^{\lft}) \To \PSh(\Rig_{S^\an}^{\lft})
\]
compatible with the respective Yoneda embeddings. Furthermore, this functor is monoidal with respect to the cartesian monoidal structure. The analogous statement for the functor $(-)^\an\colon\Sm_S\to\SmRig_{S^\an}$ holds true, too.
\end{lemma}
\begin{proof}
Restricting along the Yoneda embedding $\Sch_S^\ft \inj \Fun(\Sch_S^{\ft,\opp},\Spc)$ induces for any \infcat{} $\Ccal$ which admits small colimits an equivalence \cite[Prop. 5.1.5.6]{HTT}
\[
\Fun^\mathrm{L}\bigl( \PSh(\Sch_S^{\ft}), \Ccal \bigr) \To[\simeq] \Fun(\Sch_S^\ft,\Ccal)
\]
where $\Fun^\mathrm{L}$ denotes the full subcategory of all functors that preserve small colimits. Hence the composition $\Sch_S^\lft \To[(-)\an] \Rig_{S^\an}^\lft \inj \PSh(\Rig_{S^\an}^\lft)$ extends to an essentially unique functor 
\[
    \PSh(\Sch_S^{\lft}) \To \PSh(\Rig_{S^\an}^{\lft})
\]
via left Kan extension. Note that the analytification functor preserves finite products and hence also the left Kan extension of this functor along the Yoneda embedding.
\end{proof}

\begin{lemma} \label{analytification-homotopy-invariant--lem}
The functor from Lemma~\ref{analytification-presheaves--lem} fits into a commutative square
\carremap{\Fun(\Sch_S^{\lft,\opp},\Spc)}{}{\L_{\Arig}}{\Fun(\Rig_{S^\an}^{\lft,\opp},\Spc)}{\L_{\Arig}}{\Fun^{\Arig}(\Sch_S^{\lft,\opp},\Spc)}{}{\Fun^{\Arig}(\Rig_{S^\an}^{\lft,\opp},\Spc)}
\end{lemma}
The analogous statement for the functor $(-)^\an\colon\Sm_S\to\SmRig_{S^\an}$ holds true, too.
\begin{proof}
The functor $\Fun(\Sch_S^{\lft,\opp},\Spc) \To \Fun(\Rig_{S^{\an}}^{\lft,\opp},\Spc)$ sends a morphism $\mathrm{pr}_1\colon X\times\Arig_S\to X$ to the morphism $\mathrm{pr}_1\colon X^\an\times\Arig_{S^\an}\to X^\an$ which becomes an equivalence after applying $\L_{\Arig}$, hence we get the desired functor from the universal property of localisation. 
\end{proof}

\begin{lemma} \label{analytification-sheaves--lem}
The functor from Lemma~\ref{analytification-presheaves--lem} fits into a commutative square
\carremap{\Fun(\Sch_S^{\lft,\opp},\Spc)}{}{\L_\Nis}{\Fun(\Rig_{S^\an}^{\lft,\opp},\Spc)}{\L_\Nis}{\Fun^{\Nis}(\Sch_S^{\lft,\opp},\Spc)}{}{\Fun^{\Nis}(\Rig_{S^\an}^{\lft,\opp},\Spc)}
\end{lemma}
The analogous statement for the functor $(-)^\an\colon\Sm_S\to\SmRig_{S^\an}$ holds true, too.
\begin{proof}
The functor $\Fun(\Sch_k^{\lft,\opp},\Spc) \To \Fun(\Rig_{S^\an}^{\lft,\opp},\Spc)$ sends Nisnevich covers to Nisnevich covers so that we get the desired functor from the universal property of localisation. 
\end{proof}

\begin{lemma} \label{analytification-motivic-spaces--lem}
There exists an analytification functor $\textup{H}(S) \to \RigH(S^\an)$ fitting into a commutative square 
\carremap{\Sm_S}{(-)^\an}{}{\SmRig_{S^\an}}{}{\textup{H}(S)}{}{\RigH(S^\an).}
\end{lemma}
\begin{proof}
    This follows from Lemmas \ref{analytification-homotopy-invariant--lem} and \ref{analytification-sheaves--lem}.
\end{proof}



\begin{rem}[Change of coefficients] 
\label{analytification-coefficient-change--rem}
Let $\lambda\colon\Dcal\rightleftarrows\Ecal\colon\rho$ be an adjunction. For any category $\Ccal$, since we can compute equivalences in functor categories objectwise, we get an induced adjunction
\[ \lambda_* \,\colon\, \Fun(\Ccal\op,\Dcal) \rightleftarrows\Fun(\Ccal\op,\Ecal) \,\colon\, \rho_*.\]
If $\Ccal$ carries a topology, we get an induced adjunction
\[ \tilde{\lambda}_* \,\colon\, \Sh(\Ccal,\Dcal) \rightleftarrows\Sh(\Ccal,\Ecal) \,\colon\, \tilde{\rho}_*\]
with $\tilde{\lambda}_* \simeq \L_\Ecal\circ\lambda_*$ and $\tilde{\rho}_*\simeq\rho_*$ (omitting the inclusion functors from presheaves to sheaves), since we can patch adjunctions together. 

If $\Ccal\in\{\Sch_S^{\ft,\opp},\Sm_S,\Rig_{S^\an},\SmRig_{S^\an}\}$, then we get an induced adjunction
\[ \hat{\lambda}_* \,\colon\, \Sh^{\Arig}(\Ccal,\Dcal) \rightleftarrows\Sh^{\Arig}(\Ccal,\Ecal) \,\colon\, \hat{\rho}_*\]
on the full subcategories of $\Arig$-invariant sheaves.
Combining with Lemma~\ref{analytification-motivic-spaces--lem} we get adjunctions
\[
\textup{H}(S,\Spc) \rightleftarrows \RigH(S^\an,\Spc) \rightleftarrows \RigH(S^\an,\Pro^\omega(\Spc)).
\]
In the case, where $\Dcal$ and $\Ecal$ are commutative algebra objects in $\PrL$ and $\lambda$ is symmetric monoidal, we can give an alternative functors to the adjunction $\hat{\lambda}\rightleftarrows\hat{\rho}$, so we obtain a monoidal adjunction. We will make this precise in Section~\ref{sec.change.of.coeff.monoidal}. The main application of this alternative description will be that $\hat{\lambda}$ will be a symmetric monoidal functor.
\end{rem}



%%%%%%%%%%%%%%%%%%%%%%%%%%%%
\subsection{Functoriality}
%%%%%%%%%%%%%%%%%%%%%%%%%%%%

We have defined the motivic homotopy $\RigH(-,\Vcal)$ category with coefficients in any presentable category $\Vcal$ (Definition~\ref{RigH--definition}). However, it suffices to prove most results on functoriality for coefficients in $\Spc$.
This eases computations as we can easily reduce ourselves to the case of presheaves and localise accordingly. But it is not a drawback as we will see later in Lemma~\ref{lem-effective-coeff} that this implies that also $\RigH(-,\Vcal)$ will satisfy the same functorial properties as $\RigH$. This will include the localisation sequence proven in Theorem~\ref{thm-gluing}. 


Let $f\colon S\to T$ be a morphism of rigid spaces. Then there is a base change\footnote{Here the fibre product denotes the fibre product for rigid analytic varieties. On affinoids, this corresponds to the completed tensor product. We drop the completion from our notation.} functor 
\[ b \colon \SmRig_T\to\SmRig_S, X\mapsto X\times_TS\eqqcolon X_{S},\]
inducing by precomposition a functor $f_* := b^* \colon \PSh(\SmRig_S) \to \PSh(\SmRig_T)$ which restricts to a functor 
\[ f_* \colon \RigH(S) \To \RigH(T). \]
This functor preserves limits and is accessible\footnote{On presheaves, this is clear. Then the left adjoint of $f_{*}$ on $\RigH$ is given by $\L_{\Arig}\L_{\Nis}f^{*}$. By abuse of notation, we will again denote this left adjoint with $f^{*}$.}, hence admits a left-adjoint functor
\[ f^* \colon \RigH(T) \To \RigH(S). \]
The functor $f^*$ is a left Kan extension and for $F\in\RigH(T)$ and $X\in\SmRig_{S}$ \cite[Prop. 4.3.3.7]{HTT}, one has an equivalence
\[ (f^*F)(X) \simeq \colim_{\SmRig_T\ni Y, Y_{S}\to X}F(Y). \]
In particular, $(f^*F)(X\to S) \simeq F(X\to S\to[f] T)$ whenever $f$ is smooth. In this case the functor $f^{*}$ is induced by the restriction along the postcomposition functor $\SmRig_{T}\to \SmRig_{S}$ and as above we see that $f$ admits a left adjoint
\[ f_{\sharp}\colon \RigH(S)\to \RigH(T) \]
by formal reasons.

\begin{lemma}
\label{lem.Nis-localisation}
The category $\Sh_\Nis(\SmRig_S)$ is a topological localisation of $\PSh(\SmRig_S)$. In particular, the category $\Sh_\Nis(\SmRig_S)$ is an $\infty$-topos and the localisation functor\linebreak $\L_\Nis \colon \PSh(\SmRig_S) \to \Sh_\Nis(\SmRig_S)$ is accessible and left-exact.
\end{lemma}
\begin{proof}
This is a general fact for categories of sheaves for Grothendieck topologies, see \cite[6.2.2.7, 6.2.1.7]{HTT}.
\end{proof}

\begin{lem}
\label{lem.loc.cartesian}
	The motivic localisation $\Lmot$ preserves finite products and is locally cartesian.
\end{lem}
\begin{proof}
	For $\L_{\Nis}$ this follows from Lemma~\ref{lem.Nis-localisation} and for $\L_{\Arig}$ this follows from \cite[Prop. 3.4]{Hoy-EQ}.
\end{proof}

\begin{rem}[Monoidal structure on $\RigH$]
\label{rem-RigH-monoidal}
	Let $S$ be a rigid space. The category of spaces $\Spc$ together with its cartesian monoidal structure induces a symmetric monoidal structure on $\PSh(\SmRig_S)$ by the pointwise tensor-product \cite[Rem. 2.1.3.4]{HA}. By Lemma~\ref{lem.loc.cartesian}, the localisation $\Lmot$ is compatible with this monoidal structure, inducing a  symmetric monoidal structure on $\RigH(S)$ \cite[Prop. 2.2.1.9]{HA}. Note that the $\otimes$-product on the presheaf level commutes with arbitrary colimits and thus also on $\RigH(S)$. 
	If $f\colon S\to T$ is a map of rigid spaces, then $f^{*}$ is a symmetric monoidal functor.\footnote{This is clear as $f^{*}$ on presheaves certainly preserves finite products and as $\L_{\Arig}\L_{\Nis}$ commutes with finite products, we can reduce to this case. But this is clear, as $f^{*}$ can be written as a filtered colimit (the fact that this colimit is filtered follows from \cite[00X3]{stacks-project}).}

    Let us add that $\RigH(S)$ is obtained via localisation of an $\infty$-topos at a small set of morphisms. In particular, $\RigH(S)$ is presentable and even more it is a commutative algebra object in $\PrL$.
\end{rem}

\begin{prop} \label{prop-smooth-BC-and-PF}
	Let $f\colon S\rightarrow T$ be a smooth morphism of rigid spaces.
	\begin{enumerate}
		\item The pullback $f^{*}\colon \RigH(T)\rightarrow \RigH(S)$ admits a left adjoint $f_{\sharp}$.
		\item (Smooth projection formula) The canonical map 
			\[ f_{\sharp}(f^{*}M\otimes N)\to M\otimes f_{\sharp}N \]
			is an equivalence for all $M\in\RigH(T)$ and $N\in\RigH(S)$.
		\item (Smooth base change) Let $g\colon X\rightarrow T$ be a morphism of rigid spaces and let 
		\[ \begin{tikzcd}
			X_{S}\arrow[r,"f'"]\arrow[d,"g'"]& X\arrow[d,"g"]\\
			S\arrow[r,"f"]&T
		\end{tikzcd} \]
		be a pullback diagram. Then the exchange map $f'_{\sharp}g'^{*}\to g^{*}f_{\sharp}$ is an equivalence.
	\end{enumerate}
\end{prop}
\begin{proof}
	The first assertion follows from the discussion in the beginning. For (2) and (3) we may use that $\L_{\Arig}\L_{\Nis}$ commutes with finite products and thus, we can check them on the level of presheaves, where they follow from the construction.
\end{proof}

\begin{rem}
    Let us note that Proposition~\ref{prop-smooth-BC-and-PF} (3) is equivalent to the right adjointability of the transposed diagram, i.e. $f'_{\sharp}g'^{*}\to g'^{*}f'_{\sharp}$ is an equivalence if and only if the exchange map $f^{*}g_{*}\to g'_{*} f^{'*}$ is an equivalence.
\end{rem}

\begin{prop} \label{prop-nisnevich-descent-RigH}
The following assignments define sheaves for the Nisnevich topology on $\Rig$ with values in $\Pr^{L}$.
\begin{enumerate}
	\item The assignment $S\mapsto\Sh_\Nis(\SmRig_S), f\mapsto f^*$, and
	\item The assignment $S\mapsto\RigH(S), f\mapsto f^*.$
\end{enumerate}
\end{prop}
\begin{proof}
The first part follows from \cite[Prop. 2.3.7]{AGV}.\par 
For the second part, let us note that it is enough to show that for any Nisnevich cover $f\colon U\to X$ and any Nisnevich sheaf $F$ such that $f^*F$ is $\Arig$-invariant, also $F$ is $\Arig$-invariant. This follows since 
\[ \Homline(\Arig_{U},f^*F) \simeq f^*\Homline(\Arig_X,F) \]
and since $f^*$ is conservative, where $\Homline$ denotes the internal $\Hom$ in $\Sh_\Nis(\SmRig_S)$ (this is completely analogous to the proof of \cite[Thm. 2.3.4]{AGV}).
\end{proof}

In the rest of this subsection, we want to prove that for a closed immersion $i\colon Z\inj S$ the pushforward $i_{*}$ preserves weakly contractible colimits. This will follow via direct computation.\par
Let $S$ be a rigid space. We denote by $\es/S$ the initial object in $\SmRig_{S}$. The category of those functors $F\in \PSh(\SmRig_{S})$ such that $F(\es)\simeq \ast$ is denoted by $\PSh_{\es}(\SmRig_{S})$. Equivalently, we can localise $\PSh(\SmRig_{S})$ at the morphism $\es \to \es/S$, where $\es$ is the empty functor.  In particular, The inclusion $\PSh(\SmRig_{S})\to \PSh_{\es}(\SmRig_{S})$ admits a left adjoint, denoted by $L_{\es}$.

\begin{rem}
\label{rem.Pes.gen}
	By design\footnote{A presheaf $F\in\PSh(\SmRig_{S})$ is contained in $\PSh_{\es}(\SmRig_{S})$ if and only if its essential image is weakly contractible.}, the category $\PSh_{\es}(\SmRig_{S})$ is generated by representables under weakly contractible colimits (i.e.\ colimits of shape $I$ where the map $I\to *$ is cofinal).
\end{rem} 

\begin{lem}
	Let $i\colon Z\hookrightarrow S$ be a closed immersion of rigid spaces. The functor\linebreak $i_{\es *}\colon \PSh_{\es}(\SmRig_{Z})\rightarrow  \PSh_{\es}(\SmRig_{S})$ commutes with $\L_{\Nis}$ and $\Lmot$.
\end{lem}
\begin{proof}
That $i_{\es*}$ commutes with Nisnevich-sheafification follows with the same arguments as  \cite[Lem. 1.4.19]{AyoubRig}. For the sake of completion let us recall the proof.\par 
	It is enough to show $\L_{\Nis}\circ i_{\es*}\simeq i_{\es *}\circ \L_{\Nis}$. For this let $F\in\PSh_{\es}(\SmRig_{Z})$. Let $U$ be a rigid space over $S$. We can define
	\[
		(i_{\es*}F)^{+}(U)\coloneqq \colim_{(U_{i})_{i}\in\textup{Cov}_{\Nis}(U)}\lim_{\Delta}i_{\es*}F(\Cv(\coprod_{i} U_{i}/U)_{\bullet}),
    \]
	where the colimit runs over all Nisnevich coverings $(U_{i}\to U)_{i}$. The sheafification is then given by applying $(-)^{+}$ infinitely many times (via transfinite induction - cf. proof of \cite[Prop. 6.2.2.7]{HTT}). In the proof of \cite[Lem. 2.2.4]{AGV} it is shown that if $U\times_{S}Z$ is non-empty, then any Nisnevich covering of $U\times_{S}Z$ can be refined by the base change of a Nisnevich covering of $U$. In particular, the base change map $\textup{Cov}_{\Nis}(U)\to \textup{Cov}_{\Nis}(U\times_{S}Z)$ is cofinal. Thus, the definition of $\L_{\Nis}$ via the $(-)^{+}$-construction shows
	\[
		\L_{\Nis}(i_{\es*}F)(U)\simeq (\L_{\Nis}F)(U\times_{S}Z)\simeq i_{\es*}(\L_{\Nis}F)(U).
	\]
	\par 
	For the commutativity with $\L^{\#}_{\Arig}\circ \L_{\Nis}$ it is enough to show that for all $U\in\SmRig_{Z}$, we have that $p\colon i_{\es*}\Arig_{U}\to i_{\es*}U$ is an $\Arig_{S}$-equivalence (similar to the proof of \cite[Prop. 1.4.18]{AyoubRig}). The $0$-section of $\Arig_{U}\to U$ induces a section $s$ of $p$. We claim that also $s\circ p$ is $\Arig_{S}$-homotopy equivalent the identity on $i_{\es *}\Arig_{U}$. The composition
	\[
		\Arig_{S}\otimes i_{\es *}\Arig_{U}\to i_{\es *}\Arig_{Z}\otimes  i_{\es *}\Arig_{U}\simeq i_{\es *}(\Arig_{Z}\times_{Z}\Arig_{Z}\times_{Z} U)\xrightarrow{m}i_{\es *}(\Arig_{Z}\times_{Z} U), 
	\]
	where $m$ is the multiplication, naturally yields such a homotopy. In particular, $p$ is $\Arig_{S}$-homotopy equivalence.
\end{proof}


\begin{rem}
\label{rem.closed.no-colim}
	Let us note that $i_{\es *}$ cannot preserve the initial object, when $U\coloneqq S\setminus Z\neq \es$. This is because $i_{\es*}(\es/Z)(U)\simeq \es/Z(\es/Z)\simeq \ast \not\simeq \es = \es/S(U)$. But $i_{\es *}$ preserves weakly contractible colimits, as $\PSh_{\es}(\SmRig_{Z})$ is freely generated under weakly contractible colimits in $Z$ (and similarly for $S$).
\end{rem}

The remark above immediately implies the following.

\begin{cor} \label{cor-i_*-preserves-colimits}
Let $i\colon Z\inj S$ be a closed immersion of rigid spaces. Then the functor $i_* \colon \RigH(Z) \to \RigH(S)$ preserves weakly contractible colimits.
\end{cor}

The only obstruction for $i$ above not preserving colimits is that in $\RigH(S)$ the initial and final object may not agree. This is certainly a "problem" inside $\Spc$, which will go away after stabilising. 

%%%%%%%%%%%%%%%%%%%%%%%%%%%%%%%
\subsection{Gluing}
%%%%%%%%%%%%%%%%%%%%%%%%%%%%%%%
In this section, we want to prove the existence of a
gluing theorem in the unstable setting which yields a localisation sequence in the stable setting. Such a theorem has been first proven by Morel-Voevodsky in the algebraic case \cite[Thm.~2.21]{MV1}.  In $\Brig$-homotopy theory this result is proven by Ayoub \cite{AyoubRig}. We will follow the proof structure presented by Hoyois \cite[\S 4]{Hoy-EQ}.

\vspace{6pt}\noindent\textbf{Notation.} In this subsection, we fix a rigid space $S$. 

\begin{defi} \label{def-phi-presheaf}
Let $i\colon Z\to S$ be a closed immersion and let $U:=S\setminus Z$ be its open complement. Given a pair $(X,t)$ consisting of $X\in\SmRig_S$ and a partial section $t\colon Z\to X$, i.e.\ inducing a section $t\colon Z\to X_Z := X\times_SZ$, we define the presheaf
\[ \Phi_S(X,t) := (X\sqcup_{X_U}U) \times_{i_*X_Z} S. \]
More explicitly, evaluating at $Y\in\SmRig_S$ we have that
\begin{align*}
\Phi_S(X,t)(Y) = \begin{cases} \Hom_S(Y,X)\times_{\Hom_Z(Y_Z,X_Z)}* & (Y_Z\neq\varnothing) \\ * & (Y_Z=\varnothing) \end{cases}
\end{align*}
where $\Hom_Z(Y_Z,X_Z)$ is pointed at the map $Y_Z \to Z \to[t] X_Z$. 
\end{defi}

\begin{rem}
We get a functorial assignment $(X,t) \mapsto \Phi_S(X,t)$ and for every morphism $f\colon T\to S$ there is a natural map 
\[ f^*\Phi_S(X,t) \To \Phi_T(X_T,t_T) \]
which is an isomorphism when $f$ is smooth.
\end{rem}

\begin{lem} \label{lem-phi-presheaf-etale}
Let $p\colon X'\to X$ be an \'etale morphism in $\SmRig_S$ and let $t\colon Z\to X$ and $t'\colon Z\to X'$ be partial sections which are compatible, i.e.\ $t=p\circ t'$. Then the induced map
\[ \Phi_S(p) \colon \Phi_S(X',t') \To \Phi_S(X,t) \]
is a Nisnevich equivalence.
\end{lem}
\begin{proof}
This is \emph{Etape 2} in the proof of \cite[Prop.~1.4.21]{AyoubRig}.
\end{proof}

\begin{lem} \label{lem-phi-presheaf-contractible}
Let $n\geq 1$ and let $\sigma\colon Z \inj S \to \Anrig_S$ be the restriction of the zero section of the canonical morphism $\pi\colon\Anrig_S\to S$. Then the canonical map $\Phi_S(\Anrig,\sigma) \to S$ is an $\Arig$-equivalence.
\end{lem}
\begin{proof}
A nullhomotopy is given by the map
\[ \Arig \times \Phi_S(\Anrig,t) \To \Phi_S(\Anrig,t), \quad (a,f) \mapsto af. \]
\end{proof}


\begin{thm}[Gluing]
\label{thm-gluing}
Let $i\colon Z\inj S$ be a closed immersion of rigid spaces with open complement $j\colon U\inj S$. Then for every $F\in\RigH(S)$ the square
\carremaptag{$\heartsuit$}{j_\#j^*F}{\epsilon}{}{F}{\eta}{U}{}{i_*i^*F}
is a pushout square (where the maps without labels are the unique ones).
\end{thm}
\begin{proof}
The category $\RigH(S)$ is generated by representables under weakly contractible colimits (cf. Remark~\ref{rem.Pes.gen}) and all the functors in the square $(\heartsuit)$ preserve these colimits, either because they are left-adjoint or by Corollary~\ref{cor-i_*-preserves-colimits}. Hence we may assume that $F=\Lmot X$ for some $X\in\SmRig_S$ so that it suffices to show that the canonical map
\[ X \sqcup_{X_U}U \To i_*X_Z \]
in $\PSh(\SmRig_S)$ is motivic equivalence. 
We prove more generally that is a motivic equivalence in $\PSh(\Rig_S$).\footnote{At this point there is a mistake in the proof of  \cite[Thm.~4.18]{Hoy-EQ} because the category of smooth schemes (over a fixed base) does not have fibre products; see also the following footnote. This has been fixed in the latest version on the arXiv \cite[Thm.~4.18]{Hoy-EG-arxiv-v5}. We thank Marc Hoyois for communicating to us that the same problem appeared in a previous version of our article.}
The latter can be checked after base change along all maps $Y\to i_*X_Z$ for $Y\in\Rig_S$, i.e. it suffices to show that the maps
\[ \tag{$\spadesuit$}  (X \sqcup_{X_U}U) \times_{i_*X_Z} Y \To Y \]
are equivalences. We claim that we can reduce to the case where $Y=S$.

Indeed, given $p\colon Y\to S$ in $\SmRig_S$ and denoting by $q\colon Y_Z \to Z$ and $k\colon Y_Z\inj Y$ base changes of $p$ and $i$, respectively, we see that the map $(\spadesuit)$ identifies with the map
\[ p_\#\bigl( (p^*X \sqcup_{p^*X_U}p^*U) \times_{k_*(p^*X)_{Y_Z}} Y\bigr) \to p_\#Y \]
by a relative version of the projection formula\footnote{More precisely, for a morphism $f\colon T\to S$, presheaves $E,F\in\PSh(\Rig_S)$, and $G\in\PSh(\Rig_T)$ one has that $f_\#(f^*E\times_{f^*F}G) \simeq E\times_F G$. The proof reduces to the case of representable presheaves and then uses the existence of fibre products in $\Rig_S$.} and since we have equivalences
\[ p^*i_*X_Z \simeq k_*q^*X_Z \simeq k_*q^*i^*X \simeq k_*k^*p^*X \simeq k_*(p^*X)_{Y_Z} \]
where the first one is smooth base change (Proposition~\ref{prop-smooth-BC-and-PF}).

By adjunction, the datum of a map $S\to i_*X_Z$ corresponds to a section $t\colon Z\to X_Z$. Translating into the notation from Definition~\ref{def-phi-presheaf}, we have thus reduced to having to show that the canonical map $\Phi_S(X,t) \to S$ is a motivic equivalence for all $(X,t)$ as in loc.\! cit.

By Nisnevich descent (Proposition~\ref{prop-nisnevich-descent-RigH}) we may assume that $S$ is affine and always can replace by an open subspace. Thus we find an open neighbourhood $U$ of $Z$ in $S$, an open neighbourhood $V$ of $t(Z)$ in $X$, and an isomorphism $V\cong\Bnrig_U$ for a suitable $n\geq 1$ such that $t$ identifies with the partial zero section under this isomorphism \cite[Prop.~1.3.16]{AGV}. Applying Lemma~\ref{lem-phi-presheaf-etale} to the morphism $\Bnrig\cong V\inj X$ we obtain a Nisnevich equivalence $\Phi_S(X,t)\to[\simeq]\Phi_S(\Bnrig,\sigma)$. Applying the same lemma to the open immersion $\Bnrig\inj\Anrig$ we get a Nisnevich equivalence $\Phi_S(\Bnrig,\sigma)\to[\simeq]\Phi_S(\Anrig,\sigma)$. By Lemma~\ref{lem-phi-presheaf-contractible}, the presheaf $\Phi_S(\Anrig,\sigma)$ is $\Arig$-contractible.
\end{proof}

\begin{cor}
\label{cor-closed-im-ff}
    Let $i\colon Z\rightarrow S$ be a closed immersion. Then the functor $i_{*}\colon \RigH(Z)\to \RigH(S)$ is fully faithful. 
\end{cor}
\begin{proof}
    This is standard\footnote{Proof in other situations can be found in \cite[1.4.23]{AyoubRig}, \cite[Cor. 4.19]{Hoy-EQ} or \cite[Cor. 2.2.2 (1)]{AGV}.} and follows from Theorem~\ref{thm-gluing}, smooth base change (cf. Proposition~\ref{prop-smooth-BC-and-PF}) and Nisnevich descent. Nevertheless, let us give a proof for completion.

    
    We have to show that for any $F\in \RigH(Z)$ the counit $i^{*}i_{*}F\rightarrow F$ is an equivalence. Using Theorem~\ref{thm-gluing}, it is enough to show that $i_{*}$ is conservative. In particular, it is enough to show that for $f\colon F\rightarrow G$ in $\RigH(Z)$ such that $i_{*}f$ is an equivalence, we have that $f$ is an equivalence. By Remark~\ref{rem.Pes.gen} it is suffices to test that $f$ is an equivalence at any $X\in \SmRig_{Z}$. The proof of \cite[Lem. 2.2.5 (2)]{AGV} shows that we may assume that there exists an $S_{X}\in \SmRig_{S}$ such that $X = S_{X}\times_{S} Z$. In particular, $f(X)$ is equivalent to $i_{*}f(S_X)$.
\end{proof}

\begin{cor}
    Let us consider a cartesian square
    \[ 
    \begin{tikzcd}
			T_{Z}\arrow[r,"i'",hookrightarrow]\arrow[d,"f'"]& T\arrow[d,"f"]\\
			Z\arrow[r,"i",hookrightarrow]&S
	\end{tikzcd}
    \]
    where $i$ is a closed immersion. Then the exchange transformation $f^{*}i_{*}\to i'_{*}f'^{*}$ is an equivalence.
\end{cor}
\begin{proof}
    This follows from Corollary~\ref{cor-closed-im-ff} and the smooth base change of Proposition~\ref{prop-smooth-BC-and-PF}.
\end{proof}


\section{Stable  \texorpdfstring{$\AA^{1}$}~-homotopy theory in rigid geometry}
\localtableofcontents
\vspace{6pt}

 \noindent In this section we want to define the stable $\Arig$-homotopy category. To do so, we first want to show that the effective version satisfies all the functorial properties as $\RigH$. Afterward, we $\otimes$-invert the pointed projective line. This will yield a stable homotopy functor in the sense of Ayoub \cite{AyoubThesis}. In particular, using the results of \textit{op.\! cit.}, we will obtain a six functor formalism with respect to algebraic maps. 

\vspace{6pt}\noindent\textbf{Notation.} In this section, let $S$ be a rigid space and let $\Vcal\in\CAlgPr$. For results using the analytification functor, let $R$ be an adic ring with ideal of definition $I$ and set $B:=\Spec(R)\setminus\textup{V}(I)$ and $B^\an=\textup{Spf}(R)^\textup{rig}$.

\subsection{Change of coefficients}
\label{sec.change.of.coeff.monoidal}

In Section \ref{sec.unstable}, we constructed the rigid $\Arig$-homotopy category via $\Arig$-localisation of Nisnevich sheaves of spaces. In this short subsection, we want to explain how we can extend our results to $\Arig$-invariant Nisnevich sheaves with $\Vcal$-valued coefficients.


\begin{rem}
\label{rem-change-of-coeff}
	We have $\Sh_{\Nis}(\SmRig_{S})\otimes\Vcal\simeq \Sh_{\Nis}(\SmRig_{S},\Vcal)$ \cite[Prop. 2.4]{DrewMHM}. This is also compatible with $\Arig$-localisation, hence
	\[
        \RigH(S)\otimes\Vcal\simeq\RigH(S,\Vcal).
	\]
    Furthermore, by Remark~\ref{rem-RigH-monoidal} the category $\RigH(S)$ is a commutative algebra object in $\PrLO$. Thus, the same is true for $\RigH(X)\otimes\Vcal$. In particular, the $\otimes$-product in $\RigH(S)\otimes\Vcal$ is symmetric and commutes with arbitrary colimits in both variables.
\end{rem}

\begin{rem}
\label{rem-analyitification-monoidal}
    Let us consider the constant functor $\mathrm{const}\colon \Spc\to\Cond^{\omega}(\Spc)$. This is a map of presentable categories by Lemma~\ref{condensed-objects-presentable--lem}. As colimits in $\Cond^{\omega}(\Spc)$ can be computed pointwise, we see that $\mathrm{const}$ is colimit preserving. Moreover, $\Cond^{\omega}(\Spc)$ is the sheaf category associated to a Grothendieck topology. In particular, the localisation functor $\PSh(\ProFin^{\omega})\rightarrow \Cond^{\omega}(\Spc)$ is left exact. Therefore, $\mathrm{const}$ also preserves finite products and in particular is symmetric monoidal. Thus, the constant functor $\mathrm{const}$ is a map in $\CAlgPr$. The analytification functor of Remark~ \ref{analytification-coefficient-change--rem} yields a symmetric monoidal colimit preserving functor\footnote{We suspect that the functor $\RigH(B^{\an})\rightarrow \RigH(B^{\an},\Cond^{\omega}(\Spc))$ induced via base change with $\mathrm{const}$ is equivalent to the functor $\widehat{const}_*$ constructed Remark \ref{analytification-coefficient-change--rem}. However, we did not find a reference.} 
    $$
        \textup{H}(B,\Spc)\To[\textup{an}] \RigH(B^{\an})\To[\id\otimes\mathrm{const}] \RigH(B^{\an},\Cond^{\omega}(\Spc)).
    $$
\end{rem}

With the above discussion it is not hard to see that the base change of a pullback formalism is again a pullback formalism \cite[\S 8]{DrewMHM}. As our construction is not directly covered by \textit{loc.\!\! cit.}, let us be more precise and give a statement with proof for completion.

\begin{lem}
\label{lem-effective-coeff}
	 The induced functor 
        $$
        \RigH\otimes \Vcal \colon\Rig\op\to \CAlgPr_{/\Vcal},\quad S\mapsto \RigH(S)\otimes \Vcal,\quad f\mapsto f^{*}\otimes\id_{\Vcal}
        $$ 
        has the following properties.
	\begin{enumerate}
		\item[(1)] For any smooth morphism $f$ inside $\Rig$ the pullback $f^{*}\otimes \id_{\Vcal}$ admits a left adjoint $f_{\sharp}\otimes \id_{\Vcal}$ satisfying smooth base change and projection formula (cf. Proposition~\ref{prop-smooth-BC-and-PF}). 
		\item[(2)] For each $S\in\Rig$ and natural projection $p\colon \Arig_{S}\rightarrow X$ the functor $p^{*}\otimes \id_{\Vcal}$ is fully faithful.
	\end{enumerate}
\end{lem}
\begin{proof}
The proof is similar to \cite[Prop. 8.5]{DrewMHM} but let us give an argument nonetheless.\par
	If $f$ is smooth, then it admits a left adjoint as this is true for $\RigH$ \cite[Lem. 8.4]{DrewMHM}. The main argument is that $\RigH(S)\otimes\Vcal\simeq \Fun^{R}(\RigH(S)\op,\Vcal)$ by \cite[Prop 4.8.1.17]{HA}, where $\Fun^{R}$ denotes those functor that admit left adjoints. Composing with $f^{*,\opp}\leftrightarrows f_{\sharp}\op$ for $\RigH\op$ yields the desired adjunction. This construction allows us to interpret smooth base change and the projection formula as composing with the transpose of the exchange transformation resp. the projection formula in $\RigH$, which is an equivalence by Proposition~\ref{prop-smooth-BC-and-PF}. This shows (1) and analogously also (2). % $\Arig$-invariance.
\end{proof}




\subsection{Effective motives and localisation}

\vspace{6pt}\noindent\textbf{Notation.} 
For the rest of this section, we will assume that $\Vcal$ is furthermore stable.

\begin{defi}
	We define the \textit{effective stable homotopy category with coefficients $\Vcal$} by 
	\[
		\RigSH^{\eff}(S,\Vcal)\coloneqq\RigH(S)\otimes \Vcal.
	\]
	We denote the induced functor by $\RigSH^{\eff}_{\Vcal}$ and if $\Vcal\simeq \Sp$, then we simply write $\RigSH^\eff(S)$ and $\RigSH^{\eff}$. 
\end{defi}


\begin{rem}
	By Remark~\ref{rem-change-of-coeff}, we see that $\RigSH^{\eff}_\Vcal$ satisfies Nisnevich descent. Further, by \cite[Ex. 4.8.1.23]{HA} $\RigSH^{\eff}(S)$ is equivalent to the stabilisation of $\RigH(S)$.
\end{rem}


\begin{notation}
	Let $f\colon T\rightarrow S$ be a morphism of rigid spaces and $\Vcal$ a symmetric monoidal stable presentable category. Abusing notation, we will simply write 
    $$
    f^{*}\colon \RigSH^{\eff}(S,\Vcal)\to \RigSH^{\eff}(T,\Vcal)
    $$
    instead of $f^{*}\otimes \id_{\Vcal}$ and similarly $f_{*}$ for its right adjoint resp. $f_{\sharp}$ for its left adjoint (if it exists).
\end{notation}


\begin{rem}
	 We have seen in Remark~\ref{rem.closed.no-colim} that for closed immersions $i\colon Z\hookrightarrow S$ the pushforward $i_{*}$ on $\RigH$ does not preserve colimits in general. The obstruction is also very clear, in $\RigH$ the initial and final object do not agree. But $i_{*}$ on $\RigSH^{\eff}(S,\Vcal)$ does certainly preserve the initial object, as this obstruction vanishes. Therefore $i_{*}$ preserves colimits.
\end{rem}

\begin{notation}
	Let $i\colon Z\rightarrow S$ be a closed immersion. Then we denote the right adjoint of $i_{*}$ by $i^{!}$.
\end{notation}

From the gluing in Theorem~\ref{thm-gluing}, we immediately obtain a localisation sequence
 
\begin{prop}[Localisation]
\label{prop-localisation}
    Let $i\colon Z\inj S$ be a closed immersion with open complement $j\colon U\inj S$. Then for any $F\in\RigSH^{\eff}(S,\Vcal)$, we have fibre sequences
    \[
    \begin{tikzcd}[row sep = 0.3em, column sep = 1.7em]
        j_{\sharp}j^{*}F\arrow[r,""] & F \arrow[r,""]&i_{*}i^{*}F,\\
        i_{*}i^{!}F\arrow[r,""] &F\arrow[r,""] & j_{*}j^{*}F.
    \end{tikzcd}
    \]
\end{prop}
\begin{proof}
    The statement of the lemma is equivalent to the following statement. For any closed immersion $i\colon Z\hookrightarrow S$ in $\Rig$ with open complement $j\colon U\hookrightarrow S$, the square 
		\[
			\begin{tikzcd}
				\RigH(Z)\otimes \Vcal\arrow[r,"i_{*}\otimes\id_{\Vcal}"]\arrow[d,""]& \RigH(S)\otimes \Vcal\arrow[d,"j^{*}\otimes\id_{\Vcal}"]\\
				\ast\arrow[r,""]&\RigH(U)\otimes \Vcal
			\end{tikzcd}
		\]
		is a pullback in $\PrL$. In particular, as the functor $-\otimes \Vcal$ preserves colimits, we can use the same arguments as in the proof of Lemma~\ref{lem-effective-coeff} to assume $\Vcal\simeq \Sp$. In view of Remark~\ref{rem-change-of-coeff}, we have $\RigSH^{\eff}(-)\simeq \Sp(\RigH(-,\Spc_{*}))$, where $\Spc_{*}$ denotes the category of pointed spaces. We have an equivalence 
        $$ \Sp(\RigH(-,\Spc_{*}))\simeq \lim(\dots\xrightarrow{\Omega}\RigH(-,\Spc_{*}))
        $$
        inside $\Pr^{R}$ \cite[Prop. 1.4.2.24]{HA}. In particular, the second sequence of the proposition is a fibre sequence if and only if it is so in $\RigH(-,\Spc_{*})$. Since the first sequence is the adjoint of the second one, we may assume $\Vcal\simeq \Spc_*$.

        We will show that the first sequence is a cofibre sequence in $\RigH(-,\Spc_{*})$, which proves that the latter one is a fibre sequence. For this let us note that the forgetful functor $p\colon \Spc_*\rightarrow \Spc$ takes pushout squares to pushout squares and reflects them \cite[Prop. 4.4.2.9]{HTT}. As in the proof of \cite[Prop. 5.2]{Hoy-EQ}  let us look at the following two squares
        $$
         \begin{tikzcd}
            j_{\sharp}j^*p(F)\arrow[r,""]\arrow[d,""]&j_{\sharp}j^*p(F)\coprod_{U} S\arrow[r,""]\arrow[d,""]& p(F)\arrow[d,""]\\
				U\arrow[r,""]&S\arrow[r,""]&i_*i^*p(F).
		\end{tikzcd}
        $$
        The outer square is a pushout by Theorem~\ref{thm-gluing} (note that $F\in \RigSH^{\eff}(S,\Vcal)$ comes equipped with a zero section). The left square is by design a pushout. So the right square is also a pushout, which is also the image of the first sequence of the proposition under $p$.
\end{proof}

\subsection{Thom motives and stabilisation}
\label{sec.thom-motives-stabilisation}

 For any vector bundle $p\colon V\rightarrow S$ with zero section $s\colon S\rightarrow V$, we can define the suspension $\Sigma_{V}\coloneqq p_{\sharp}s_{*}$ and loop $\Omega_{V}\coloneqq s^{!}p^{*}$. Note that this yields an adjunction
\[
	 \begin{tikzcd}
		\Sigma_{V}\colon\RigSH^{\eff}(S,\Vcal)\arrow[r,"",shift left = 0.3em]&\arrow[l,"",shift left = 0.3em]\RigSH^{\eff}(S,\Vcal)\colon \Omega_{V}.
	\end{tikzcd}
\]  

\begin{lem}
\label{lem-thom}
     Then the following are equivalent.
    \begin{enumerate}
        \item[(i)] For any smooth morphism $f\colon T\rightarrow S$ of rigid spaces that admits a section $s$, the Thom transformation $\textup{Th}(f,s)\coloneqq f_{\sharp}s_{*}$ is an equivalence.
        \item[(ii)] For any vector bundle $V$ on $S$ the suspension of the unit $\Sigma_{V}1_V$ is $\otimes$-invertible inside $\RigSH^{\eff}(S,\Vcal)$.
        \item[(iii)] The suspension of the unit of the closed unit disc $\Sigma_{\BB^{1}_S}1_{\BB^{1}_{S}}$ is $\otimes$-invertible.
       % \item[(iv)] The Thom transformation $\textup{Th}(f,s)$ with respect to the rigid affine line $f\colon \Arig_{S}\rightarrow S$ and the zero section $s$ is an equivalence.
    \end{enumerate}
\end{lem}
\begin{proof}
    Equivalences $(i)\Leftrightarrow (ii)\Leftrightarrow (iii)$ are standard and follow from the fact that any smooth morphism $f\colon T\rightarrow S$ of rigid spaces factors locally as an \'etale morphism $T\rightarrow \BB^{n}_S$ followed by the projection \cite[\S 2.4]{CD1}.
    %Now note that in particular any smooth morphism factors locally as an \'etale morphism $T\rightarrow \Arig_{S}$ followed by the projection. Thus, by the same argumentations $(i)\Leftrightarrow (iv)$.
\end{proof}


\begin{rem}
    Let $V\rightarrow S$ be a vector bundle. The localisation sequence in Proposition~\ref{prop-localisation} implies that 
    $$
    \Sigma_{V}1_{V}\simeq \frac{V}{V\setminus \lbrace 0\rbrace},
    $$
    where the RHS denotes the cofibre of the natural open immersion $V\setminus\lbrace 0\rbrace \rightarrow V$ induced by the zero section.
    
    Further, note that we have a commutative diagram with pushout squares
     $$
         \begin{tikzcd}
            \Brig_{S}\setminus \lbrace 0\rbrace\arrow[r,""]\arrow[d,""]&\Arig_{S}\setminus \lbrace 0\rbrace\arrow[r,""]\arrow[d,""]& 0\arrow[d,""]\\
				\Brig_{S}\arrow[r,""]&\Arig_{S}\arrow[r,""]&\frac{\Arig_{S}}{\Arig_{S}\setminus \lbrace 0\rbrace}.
		\end{tikzcd}
        $$
        In particular, we have
        $$
           \frac{\Brig_{S}}{\Brig_{S}\setminus \lbrace 0\rbrace}\simeq \frac{\Arig_{S}}{\Arig_{S}\setminus \lbrace 0\rbrace}
        $$
        inside $\RigSH^{\eff}(S)$. Thus, condition (iii) in Lemma \ref{lem-thom} is equivalent to $\Sigma_{\Arig_S}1_{\Arig_S}$ being $\otimes$-invertible.

        Finally, let us look at the diagram
         $$
         \begin{tikzcd}
            \Arig_{S}\setminus \lbrace 0\rbrace\arrow[r,""]\arrow[d,""]&\Arig_{S}\arrow[r,""]\arrow[d,"\infty"]& 0\arrow[d,""]\\
				\Arig_{S}\arrow[r,"0"]&\PP^{1}_{S}\arrow[r,""]&(\PP^{1}_{S},\infty),
		\end{tikzcd}
        $$
        where $(\PP^{1},\infty)$ is defined as the pushout of the right square. All squares in the diagram are a pushout thus we have 
        $$
            (\PP^{1}_{S},\infty)\simeq  \frac{\Arig_{S}}{\Arig_{S}\setminus \lbrace 0\rbrace}.
        $$
\end{rem}


\begin{defi}
	We define the category of \textit{rigid analytic motivic spectra over $S$}
	\[
		\RigSH(S,\Vcal)\coloneqq\RigSH^{\eff}(S,\Vcal)[(\PP^{1}_{S},\infty)^{-1}].
	\]
\end{defi}



\begin{rem}
\label{rem-RigSH-mon}
    As remarked in \cite[\S 6.1]{Hoy-EQ}, the definition of inversion along an element yields for a morphism of rigid spaces $f\colon S\rightarrow T$ an equivalence
    $$
        \RigSH^{\eff}(S,\Vcal) \otimes_{\RigSH^{\eff}(S,\Vcal)}\RigSH(T,\Vcal)\simeq \RigSH(S,\Vcal)
    $$
    via $f^*$ as symmetric monoidal categories. In particular, the argumentation of the proof of Lemma \ref{lem-effective-coeff} shows that $\RigSH(-,\Vcal)$ again satisfies projection formula, smooth base change formula and localisation. We refer to \cite[Prop. 8.5]{DrewMHM} for a rigorous proof of this statement. Furthermore, by construction $(\PP^{1}_{S},\infty)$ is $\otimes$-invertible inside $\RigSH(S,\Vcal)$.

    Lastly, let us remark that by design $\RigSH(S,\Vcal)$ is a commutative $\Vcal$-algebra in $\PrL$ \cite[Def. 2.6]{Robalo-Thesis} as $\RigSH^{\eff}(S,\Vcal)\in \CAlgPr_{/\Vcal}$ by Remark \ref{rem-change-of-coeff}.
\end{rem}



Using this definition, we can conclude that $X\mapsto \RigSH(X^{\an})$, for any $B$-scheme $X$, is a stable homotopy functor \cite{AyoubThesis}  (see also \cite[Prop. 5.11]{DrewMHM}), i.e. it satisfies the properties (1)-(6) listed in \cite[\S 1.4.1]{AyoubThesis}.\par
In particular, \cite[Scholie 1.4.2]{AyoubThesis} implies the existence of a six functor formalism on the restriction of $\RigSH$ to quasi-projective $K$-schemes. Using the results of Cisinksi-Deglise, we can extend this to all separated finite type $B$-schemes \cite[Thm. 2.4.50]{CD1}.

\begin{thm}
\label{thm-6ff}
The functor\footnote{The functoriality follows similarly as in \cite[\S 9.1]{RobaloDiss}.} 
\begin{align*}
	\RigSH_{\Vcal}\colon (\Sch^{\textup{sep,ft}}_{B})\op\rightarrow \CAlgPr_{/\Vcal},\quad X\mapsto \RigSH(X^{\an},\Vcal),\quad f\mapsto f^{*}
\end{align*}
satisfies Nisnevich descent and extends to a full six functor formalism in the sense of \cite{Liu-Zheng}. Further, by passing to homotopy categories this functor defines a motivic triangulated functor in the sense of \cite{CD1}.
\end{thm}
\begin{proof}
    Nisnevich descent follows from Proposition \ref{prop-nisnevich-descent-RigH}.

    By \cite[Prop. 5.11]{DrewMHM} together with Remark \ref{rem-RigH-monoidal}, Lemma \ref{lem-effective-coeff}, Proposition \ref{prop-localisation} and Lemma \ref{lem-thom} the composition of $\RigSH_{\Vcal}$ with the homotopy functor yields a stable homotopy $2$-functor in the sense of \cite[\S 1.4.1]{AyoubThesis}. Thus, using that any finite type separated $B$-scheme admits a Nagata compactification \cite[Thm. 4.1]{NagataComp}, we can apply \cite[Thm. 2.4.26]{CD1} and \cite[Prop A.5.10]{Mann-Thesis} to obtain a pre-six functor formalism in the sense of \cite{Mann-Thesis}. To lift this to a full six functor formalism, we are left to show that for any proper map of separated finite type $B$-schemes $f\colon X\rightarrow Y$ the functor $f_*$ admits a right adjoint \cite[Prop. A.5.10]{Mann-Thesis}. 
    
    Let us assume that $\Vcal\simeq \Sp$, then this follows from \cite[Prop. 5.17]{DrewMHM}. Now let $\Wcal$ be a symmetric monoidal stable presentable category. Since $f_{*}$ is a morphism inside $\Pr^{L}$, so is the right adjoint $f_{*}\otimes \id_{\Wcal}$ of $f^{*}\otimes \id_{\Wcal}$ \cite[Lem. 8.4]{DrewMHM}. In particular, $f_{*}\otimes \id_{\Wcal}$ admits a right adjoint.
\end{proof}

\subsection{The relation with $\PP^1$-spectra}
\label{sec-RigSH-as-spectra}

In this subsection, we shortly want to relate our definition of the stable motivic homotopy category with the more classical approach via $\PP^1$-spectra. This will be an easy consequence from our constructions and the analytification functor
$$
    (-)^{\an}\colon \textup{SH}^{\eff}(B)\to \RigSH^{\eff}(B^\an).
$$

Let us first recall what we mean with $\PP^1$-spectra following Hovey \cite{Hovey}.\linebreak A slightly more thorough discussion can be found in \cite[\S 2.3]{RobaloDiss}. Let $\Ccal$ be a combinatorial simplicial symmetric monoidal model category, i.e. a presentable symmetric monoidal category. For any (cofibrant) object $X\in\Ccal$, in fact for any left Quillen endofunctor, one can define the model category of $X$-spectra in $\Ccal$, denoted by $\Sp^{\NN}(\Ccal, X)$. The underlying $1$-category consists of \textit{$X$-spectra} $F$, which is a sequence $(F_0,F_1,\dots)$ in $\Ccal$ together with morphisms $X\otimes F_n\rightarrow F_{n+1}$ for all $n\in \NN_0$. The notion of morphism of $X$-spectra is the obvious one. We endow $\Sp^{\NN}(\Ccal,X)$ with the stable model structure \cite[\S 3]{Hovey}. Let us remark that in general $\Sp^{\NN}(\Ccal,X)$ is not symmetric monoidal. For this reason one usually passes to \textit{symmetric spectra} together with the stable model structure, denoted by $\Sp^\Sigma(\Ccal,X)$. We do not want to go more into detail concerning $\Sp^\Sigma(\Ccal,X)$ and refer to \cite[\S 7]{Hovey}. It is important to remark that the associated categories via the simplicial nerve of $\Sp^\Sigma(\Ccal,X)$ and $\Sp^\NN(\Ccal,X)$ are not equivalent in general. However, as discussed in \cite[\S 10]{Hovey}, one can see that the only obstruction is the cyclic permutation $(123)\in \Sigma_3$ self map of $X^{\otimes 3}\coloneqq X\otimes X\otimes X$. To be more precise, Hovey shows that there is a zig-zag of Quillen equivalences between $\Sp^\Sigma(\Ccal,X)$ and $\Sp^\NN(\Ccal,X)$ if the cyclic permutation on $X\otimes X\otimes X$ is homotopic to the identity. In this case $X$ is called \textit{symmetric}. Moreover, Robalo shows that if $X$ is symmetric, we have an equivalence 
$$
    \Ccal[X^{-1}]\To[\simeq] \Sp^\Sigma(\Ccal,X)
$$
as presentable symmetric monoidal categories \cite[Thm. 2.26]{Robalo-Thesis}.

Now let us come back to our situation and show how we can interpret $\RigSH$ as $\PP^1$-spectra. For this consider the category $\RigH(B^\an)^{\wedge}_*\coloneqq \RigH(B^\an,\Spc_*)$ with the induced $\otimes$-product via Remark \ref{rem-change-of-coeff}. We also denote by $\Gm\coloneqq \Arig\setminus\lbrace 0\rbrace$ the analytic affine line without $0$.

\begin{lem}
\label{lem-P1-spectrum}
    The pointed projective line $(\PP^1,\infty)\in \RigH(B^\an)^{\wedge}_*$ is symmetric and equivalent to $S^1\otimes (\Gm,1)$.

    In particular, we have 
    $$
        \RigSH(B^{\an})\simeq \Sp^\NN(\RigH(B^\an)^{\wedge}_*,(\PP^1,\infty)).
    $$
\end{lem}
\begin{proof}
    Combining our discussion in Section \ref{sec.change.of.coeff.monoidal} and Lemma \ref{analytification-motivic-spaces--lem}, we see that  $(-)^{\an}$ is a monoidal functor. The cyclic permutation of $(\PP^1,\infty)^{\otimes 3}\in  \textup{H}(B)^\wedge_*$ is equivalent to the identity \cite[Lem. 4.4]{voevodsky-icm-1998}. In particular, the same is true for $(\PP^{1},\infty)^{\an}$. As the analytification functor preserves colimits, we see that $(\PP^{1},\infty)^{\an}\simeq (\PP^1,\infty)\in \RigH(B^\an)^{\wedge}_*$.

    Moreover, by the same argument we have 
    $$
        (\PP^1,\infty)\simeq S^1\otimes(\Gm,1)\in \RigH(B^\an)^{\wedge}_*
    $$
    \cite[Lem. 2.15]{MV1}. Therefore, we have 
$$
    \RigSH(B^\an)\simeq \RigH(B^\an)^{\wedge}_*[(\PP^1,\infty)^{-1}],
 $$
as presentable symmetric monoidal categories.
\end{proof}

\begin{cor}
    There is an equivalence of presentable symmetric monoidal categories 
    $$
        \RigSH(B^{\an},\Vcal)\simeq \Sp^\NN(\RigSH^{\eff}(B^\an,\Vcal),(\PP^1,\infty)).
    $$
\end{cor}
\begin{proof}
    This follows from Lemma \ref{lem-P1-spectrum}, the monoidality in Section \ref{sec.change.of.coeff.monoidal} and the discussion in the beginning of this subsection.
\end{proof}

 



%%%%%%%%%%%%%%%%%%%%%%%%%%%%%%%%%%%%%%%%%%%%%%%%%
\section{Applications to the K-theory of rigid spaces}
%%%%%%%%%%%%%%%%%%%%%%%%%%%%%%%%%%%%%%%%%%%%%%%%%
\localtableofcontents
\vspace{6pt}\noindent\textbf{Notation.}
In this entire section, let $(R,R^+)$ be a Huber pair such that $R$ is a Tate ring admitting a ring of definition $R_0$ which is noetherian and of finite dimension.\footnote{We need this assumption since pro-cdh descent is known only for noetherian schemes of finite dimension. Once the latter is known in a greater generality, our results will be known as well in this greater generality.}
We set $S=\Spa(R,R^+)$ and choose a pseudo-uniformiser $\pi\in R$ so that $R=R_0[\pi^{-1}]$.

%%%%%%%%%%%%%%%%%%%%%%%%%%%%%%%%%%%%%%%%%%%%%%%%%%%%%%
\subsection{Continuous K-theory and analytic K-theory}
\label{subsec.continuous-analytic-k-theory}
%%%%%%%%%%%%%%%%%%%%%%%%%%%%%%%%%%%%%%%%%%%%%%%%%%%%%%
In this section we recall the definition and basic properties of analytic K-theory for rigid spaces due to Kerz-Saito-Tamme \cite{kst-i,kst-ii}. 
Some results use the following assumption, cf.\ \cite[\S 3.2]{kst-i}.

\vspace{6pt}\noindent\textbf{Assumption $(\dagger)_R$.}
\label{assumption-dagger}
Every regular and topologically of finite type $R$-algebra $A$ admits a noetherian ring of definition $A_0$ and an admissible desingularisation, i.e.\ a proper morphism of schemes $p\colon X\to\Spec(A_0)$ with $X$ regular such that the restriction $p^{-1}(\Spec(A))\to\Spec(A)$ is an isomorphism.

%%%%%%%%%%%%%%%%%%%%%%%%%%
\subsubsection*{Pro-categories} \label{pro-categories--subsubsection}
%%%%%%%%%%%%%%%%%%%%%%%%%%



Analytic K-theory of rigid spaces is naturally a pro-object in the category of spectra. Let us quickly recall some facts about pro-categories  \cite[\S A.8.1]{SAG}. Let $\Ccal$ be a presentable category. Then we can define $\Pro(\Ccal)$ as those functors in $\Fun(\Ccal,\Spc)$, which are accessible and preserve finite limits. Equivalently, we have $\Pro(\Ccal)\simeq \textup{Ind}(\Ccal\op)\op$. An object $X\in\Pro(\Ccal)$ can be thought of a limit of a filtered  diagram $\lbrace X_{i}\rbrace$, where each $X_{i}$ is contained in the essential image of the Yoneda embedding $\Ccal\rightarrow \Pro(\Ccal)$. Following \cite{kst-i}, we denote these objects as $\prolim X_{i}$.
We have a limit functor $\Pro(\Ccal)\rightarrow \Ccal$ that maps an object $\prolim X_i$ to $\lim X_i \in \Ccal$. 

Let us remark, that finite limits and colimits in $\Pro(\Ccal)$ can be computed pointwise in the following sense. Let $K$ be a finite simplicial set. The limit functor $\Fun(K,\Ccal)\to \Ccal$ induces a functor $\Pro(\Fun(K,\Ccal))\rightarrow \Pro(\Ccal)$, which we call \textit{level-wise limit} (similarly for colimits). Furthermore, there exists a functor $\Pro(\Fun(K,\Ccal))\rightarrow \Fun(K,\Pro(\Ccal))$. Using this functor we obtain that the levelwise limit of $F\in \Pro(\Fun(K,\Ccal))$ is the limit of the induced diagram $K\rightarrow \Pro(\Ccal)$ (similarly for colimits) \cite[Lem. 2.1]{kst-i}.

Let $\Ccal = \Sp$. Then $\Pro(\Sp)$ is stable (Lemma~\ref{Pro-omega-stable--lemma}) and we can define homotopy groups on $\Pro(\Sp)$ pointwise \cite{kst-i}. To be more precise, for any $n\in \ZZ$ and any $\prolim X_i\in \Pro(\Sp)$, we set 
$$
    \pi_{n} (\prolim X_i )\coloneqq \prolim \pi_{n} X_{i} \in \Pro(\textup{Ab}).
$$
We will see that $\Pro(\Sp)$ admits a $t$-structure  and the pointwise homotopy groups are precisely the homotopy groups induced by this $t$-structure (Lemma \ref{lem-t-structure-pro}).

Lastly, let us recall the notion of \textit{weak equivalence}. In \textit{op.\!\! cit.} it is shown that there is a functor $\iota^{*}\colon \Pro(\Sp)\rightarrow \Pro(\Sp^{+})$, where $\Sp^{+}$ denotes the category of bounded above spectra. On objects this functor is given by 
$$
    \iota^{*}(\prolim_I X_i)\simeq \prolim_{\NN\times I} \tau_{\leq n}X_{i}.
$$
A morphism of pro-spectra $f\colon X\rightarrow Y$ is a weak equivalence if $\iota^{*}f$ is an equivalence. Analogously, a fibre sequence in $\Pro(\Sp)$ is called a \textit{weak fibre sequence} if it is a fibre sequence after applying $\iota^{*}$.
Note that $f$ is  weak equivalence if and only if $\tau_{\leq n}f$ is an equivalence for some $n\in\ZZ$ and $\pi_{i}f$ is an isomorphism for all $i\in \ZZ$ \cite[Lem. 2.8]{kst-i}.

\begin{defi} \label{Def:Pro-omega}
For an \infcat{} $\Ccal$ denote by $\Pro^\omega(\Ccal)$ the full subcategory of $\Pro(\Ccal)$ of objects $X$ that are representable by diagrams $X \colon I\to\Ccal$ for a \emph{countable} cofiltered category $I$.
\end{defi}

\begin{lemma} \label{Pro-omega-stable--lemma}
Let $\Ccal$ be a stable category. Then the categories $\Pro(\Ccal)$ and $\Pro^\omega(\Ccal)$ both are stable.
In particular, the categories $\Pro(\Sp^+)$ and $\Pro^\omega(\Sp^+)$ are stable.
\end{lemma}
\begin{proof}
Since $\Ccal$ is stable, the pro-category $\Pro(\Ccal)$ is stable as well \cite[2.5]{kst-i}. Since finite colimits in pro-categories can be computed levelwise, the inclusion $\Pro^\omega(\Ccal) \inj \Pro(\Ccal)$ creates finite colimits. Hence $\Pro^\omega(\Ccal)$ has admits finite colimits and the suspension functor is an equivalence. Thus $\Pro^\omega(\Ccal)$ is stable \cite[1.4.2.27]{HA}.
The rest follows as the category $\Sp^+$ is a stable subcategory of $\Sp$ \cite[p.~44]{HA}.
\end{proof}

In the work of Kerz-Saito-Tamme, the notion of \emph{pro-homotopy invariance} plays a major role when working with rings \cite{kst-bass-quillen}, \cite[\S 5.3]{kst-i}. For sheaves, pro-homotopy invariance is equivalent to $\Arig$-invariance.

\begin{lemma} \label{pro-invariant=Aan-invariant--lem}
A sheaf $F \colon \Rig_S\op \to \Pro(\Ccal)$ is $\Arig$-invariant if and only if it is pro-$\Brig$-invariant, i.e. if the canonical map 
    \[
	F(X) \To \prolim_{t\mapsto\pi t} F(X\times\Brig)
	\]
is an equivalence in $\Pro(\Ccal)$.
\end{lemma}
\begin{proof}
As in the proof of Lemma~\ref{B1-implies-Aan-invariance--lem}, we compute\footnote{Note that a pro-system $\prolim_iC_i$ is the limit of the constant pro-systems given by the $C_i$ in $\Pro(\Ccal)$.}
	\begin{align*}
	F(X\times\Arig) 
	\simeq F(\colim_{t\mapsto\pi t}(X\times\Brig)) 
	\simeq \prolim_{t\mapsto\pi t} F(X\times\Brig).
	\end{align*}
Hence the map $F(X) \to F(X\times\Arig)$ is an equivalence if and only if the map $F(X) \to \prolim_{t\mapsto\pi t} F(X\times\Brig)$ is an equivalence.
\end{proof}



\subsubsection*{Connective analytic $\K$-theory}
The idea of (connective) analytic $\K$-theory is to define an analogue of topological K-theory for rigid spaces. The new idea in \cite{kst-i} is to view analytic $\K$-theory as an $\Arig$-invariant object. This construction can be enhanced to endow analytic $\K$-theory with the structure of a pro-spectrum via (\ref{eq-Arig-colim}). Let us be more precise:
We define the \textit{standard $n$-simplex with radius $1/\pi^{m}$} as
$$
    \Delta^{n}_{m} \coloneqq \Spa(R\langle \pi^{m}T_{0},\dots, \pi^{m}T_{n}\rangle/(T_0+\dots + T_n -1), R^+\langle \pi^{m}T_{0},\dots, \pi^{m}T_{n}\rangle/(T_0+\dots + T_n -1)).
$$
There exists a (non-canonical) $S$-isomorphism between $\Delta^{n}_{m}$ and $\BB^{n}_{1/\pi^m}$.
Via the standard structure maps we obtain a cosimplicial rigid space $\Delta^\bullet_m$.
Let $X$ be a rigid space. The functor $\Delta\op\to \Sp_{\geq 0}$ given by $[n]\mapsto \textup{K}_{\geq 0}(X\times \Delta^{n}_{m})$, where $\textup{K}_{\geq 0}(-)$ denotes the connective algebraic K-theory spectrum,  defines via geometric realisation a connective spectrum 
$$
    \K_{\geq 0}(X\times \Delta_{m})\coloneqq \colim_{\Delta\op} \K_{\geq 0}(\Ocal(X\times \Delta^{\bullet}_{m})).
$$
Taking the limit over $m\geq 0$, we can define connective analytic $\K$-theory as the pro-spectrum
$$
    \kan(X)\coloneqq \prolim_{m} \K_{\geq 0} (X\times \Delta_{m}).
$$
By design, connective analytic $K$-theory is pro-$\Brig$-invariant \cite[Prop. 6.3]{kst-i} and hence $\Arig$-invariant (Lemma \ref{pro-invariant=Aan-invariant--lem}), i.e. 
$$ \kan(X)\simeq  \kan(X\times \Arig). $$

Applying the same construction to the connected cover $\K_{\geq 1}$ we obtain the pro-spectrum $\kan_{\geq 1}(X)$. Note however that $\kan_{\geq 1}(X)$ is not the connected cover of $\kan_{\geq 0}(X)$. Nevertheless, we obtain a weak fibre sequence
\begin{align}
\label{fibre-sequence-1-connective-cover}
    \kan_{\geq 1}(X)\rightarrow \kan(X)\rightarrow \K_{0}(\Ocal(X)).
\end{align} 


\subsubsection*{Analytic K-theory}
Non-connective analytic K-theory is defined by delooping 1-connective analytic K-theory via the analytic Bass construction \cite[\S 6.3]{kst-i}.
% So geht's :) thanks :) Aber habe eine Frage auf Telegram^^ Gesehen
Let us denote by $\BB^{-1}$ the adic space $\Spa(R\langle t^{-1}\rangle,R^+\langle t^{-1}\rangle)$. We define $\Lambda\kan_{\geq 1}$ via the following fibre sequence of pro-spectra  
$$
    \Lambda \kan_{\geq 1}(X)\rightarrow \kan_{\geq 1}(X \times \Brig)\sqcup_{\kan_{\geq 1}(X)} \kan_{\geq 1}(X \times \BB^{-1})\rightarrow \kan_{\geq 1}(X \times \Brig \times \BB^{-1}).
$$
Iterating this process, we can define (non-connective) analytic K-theory via the analytic Bass construction 
$$
    \Kan(X)\simeq \colim (\kan_{\geq 1}(X)\xrightarrow{\lambda} \Lambda \kan_{\geq 1}(X)\xrightarrow{\Lambda(\lambda)}  \Lambda^2 \kan_{\geq 1}(X)\rightarrow\dots),
$$
for more details see section~\ref{sec:analytic-bass-construction} below. 
Though not a purely formal consequence, it turns out that analytic K-theory satisfies the Bass Fundamental Theorem, i.e. for any $i\in\ZZ$ and any affinoid rigid space $X = \Spa(A,A^+)$ over $S$ we have an exact sequence
$$
    0\rightarrow \Kan_{i}(X)\rightarrow \Kan_{i}(X \times \Brig)\oplus \Kan_{i}(X \times \BB^{-1})\rightarrow \Kan_{i}(X \times \Brig \times \BB^{-1})\rightarrow \Kan_{i-1}(X)\rightarrow 0
$$
and the map $\Kan_{i}(X \times \Brig \times \BB^{-1})\rightarrow \Kan_{i-1}(X)$ admits a split \cite[Cor.~2.6]{kst-ii}. 
%In particular, we get an isomorphism $\lambda \colon \Kan_{i}(X)\rightarrow \Sigma \Kan_{i+1}(X)$ of pro-groups.



Again, we define the \textit{$i$-th analytic $\K$-group of $X$} as 
\[ \Kan_{i}(X) := \pi_{i}\Kan(X)\in \Pro(\textup{Ab}). \]
%Note that by design, we have $\Kan_{i}\simeq \pi_i\kan_{\geq 1}$ for $i\geq 1$.

Alternatively, we could apply the Bass construction to $\kan(X)$ and get a pro\hyp{}spectrum $(\kan(X))^{B}$. In general is not clear if both constructions are equivalent. However, if $X$ is regular and we assume $(\dagger)$, then $\pi_{i} (\kan(X))^{B} \cong \Kan_{i}(X)$ for all $i\geq 0$. We will recall below from \cite{kst-ii} that the construction via $\kan_{\geq 1}$ yields better formal properties.

\subsubsection*{Continuous $\K$-theory}
The definition of continuous $\K$-theory is due to Morrow \cite{Morrow} and the theory has been elaborated by Kerz-Saito-Tamme \cite{kst-i}.

Let $A_0$ be $\pi$-adic ring for some $\pi\in A_0$.  Then we set $\Kcont(A_0)\coloneqq \prolim \K(A_0/\pi^{n})$ which is an object in $\Pro^\omega(\Sp)$. For a Tate ring $A$ with ring of definition $A_0$, we can define $\Kcont(A)$ as the pushout of the diagram , up to weak equivalence
$$ \begin{tikzcd}
		\K(A^{\circ})\arrow[r,""]\arrow[d,""]& \K(A)\arrow[d,""]\\
		\Kcont(A_0)\arrow[r,""]&\Kcont(A).
\end{tikzcd} $$
within the category $\Pro^\omega(\Sp)$. When passing to the localisation $\Pro^\omega(\Sp^+)$, the object $\Kcont(A)$ does not depend on the choice of $A_0$ \cite[Prop.~5.4]{kst-i}.


\begin{rem}(K-theory of nuclear modules) \label{connection-efimov-k-theory--rem}
As a matter of fact, the underlying spectrum $\underline{\K}^\textup{cont}(A)$, i.e.\ the image under the limit functor $\Pro^\omega(\Sp)\to\Sp$, identifies with the K-theory of the dualisable category $\textup{Nuc}(A_0)$ of nuclear $A_0$-modules. As a consequence, we further get an equivalence $\underline{\K}^\textup{cont}(A)\simeq \K(\textup{Nuc}(A))$. For more details, we refer to Andreychev's thesis \cite[Satz~5.8]{andreychev-thesis}; unfortunately, there seems to be no full proof in the literature.
\end{rem}


\begin{prop}
    If $X$ is regular and assuming $(\dagger)_R$, then $\Kcont$ is $\Arig$-invariant, i.e.\ the canonical map $\Kcont(X) \to[\simeq] \Kcont(X\times \Arig)$ is an equivalence in $\Pro(\Sp^+)$. Moreover, the connective cover of $\Kcont(X)$ agrees with $\kan(X)$.
\end{prop}
\begin{proof}
The claims can be checked locally. In this case, $\Kcont(X)\simeq\Kan(X)$ \cite[Thm.~6.19]{kst-i} and $\Kan$ is pro-$\Brig$-invariant by design, hence $\Arig$-invariant by Lemma~\ref{pro-invariant=Aan-invariant--lem}.
\end{proof}





\subsubsection*{Descent for analytic $K$-theory}

\begin{lem}[Nisnevich descent for analytic K-theory]
The presheaf
\[ \Kan \,\colon\, \Rig_S^{\lft,\opp} \To \Pro(\Sp^+),\quad \Spa(A,A^+)\mapsto\Kan(A), \]
is a Nisnevich sheaf.
\end{lem}
\begin{proof}
Kerz-Saito-Tamme have shown that analytic K-theory on affinoid adic spaces over $S$ satisfies Zariski-descent and extends essentially uniquely to a Zariski sheaf on all adic spaces that are locally of finite type over $S$ \cite[Thm. 4.1]{kst-ii}.

For Nisnevich descent, we note that every object in $\Rig_S^\lft$ is an analytic adic space. Thus we can check whether a presheaf is a Nisnevich sheaf on elementary Nisnevich squares \cite[Kor.~A.28]{andreychev-thesis}. 

Let $(A,A^{\circ})$ be a Huber pair such that $A$ is Tate. Recall, that we assume that $A^{\circ}$ is noetherian and satisfies ($\dagger$). Then there is a fibre sequence \cite[Thm. 2.10]{kst-ii}
$$
	\K^{\pi^{\infty}}(A)\rightarrow \Kcont(A)\rightarrow \Kan(A)
$$
where 
$$
	\K^{\pi^{\infty}}(A)\coloneqq \prolim_{m}\colim_{[n]\in \Delta}\K(A\otimes_{R} R[T_{0},\dots,T_{n}]/(T_{0}+\dots+T_{n} - \pi^{m})).
$$
Hence one can reduce Nisnevich descent of $\Kan$ to Nisnevich descent of $\K^{\pi^{\infty}}$ and $\Kcont$. Both $\K^{\pi^{\infty}}$ and $\Kcont$ are induced by pro-systems of algebraic K-theory. As the \v{C}ech nerve of an open cover is $n$-skeletal, we may use that finite limits in $\Pro(\Sp^+)$ can be computed levelwise and reduced to Nisnevich descent of algebraic K-theory.
\end{proof}


%%%%%%%%%%%%%%%%%%%%%%%%%%%%%%%%%%%%%
\subsection{Identifying the representing object of analytic K-theory}
\label{subsec:representability}
%%%%%%%%%%%%%%%%%%%%%%%%%%%%%%%%%%%%%

Analytic K-theory $\Kan\colon\Rig_S^{\lft,\opp}\to\Pro^\omega(\Sp^+)$ is an $\Arig$-invariant Nisnevich sheaf, see the preceding subsection. Since both $\tau_{\geq 1}$ and $\Omega^\infty$ are right adjoint functors, the induced functor 
\[ \Rig_S^{\lft,\opp}\to\Pro^\omega(\Sp^+) \To \Pro^\omega(\Sp) \To[\tau_{\geq 1}] \Pro^\omega(\Sp_{\geq 1}) \To[\Omega^\infty] \Pro^\omega(\Spc) \]
is again an $\Arig$-invariant Nisnevich sheaf, hence an object in $\RigH(k,\Pro(\Spc))$ which we denote by $\Omega^\infty\Kan_{\geq 0}$. We would like to identify this object with more concrete objects. Unfortunately, we only can do this under the assumption $(\dagger)_{R}$ in the beginning of Section~\ref{assumption-dagger}. In this case we have for any $R$-algebra $A$ of topologically finite type that \cite[Lem.~7.5]{kst-i}
\[
\Omega^\infty\tau_{\geq 1}\Kan(A) \simeq \KVan(A) := \prolim_\rho\, \BGL(A\skp{\Delta}_\rho).
\]
This equivalence enables us to identify analytic K-theory similarly as the classical identification of algebraic K-theory with $\ZZ\times\BGL$ in the Morel-Voevodsky category \cite[Prop.~3.10]{MV1}.


\begin{thm}[Representability]
\label{thm:representability}
We assume $(\dagger)_R$.
%\hlcan{Brauchen wir $(\dagger)$? Das Setting von KST-2 sollte doch ausreichend sein.} 
Then there is a canonical equivalence
\[ \Omega^\infty\Kan_{\geq 0} \simeq \Lmot(\ZZ\times\BGL) \]
in the category $\RigH(S)$.
In particular, for every rigid space $X$ there is a functorial equivalence
\[ \Omega^\infty\Kan_{\geq 0}(X) \simeq \Homline_{\Cond^\omega(\Spc)}(\Lmot X,\Lmot(\ZZ\times\BGL)) \]
in the category $\Cond^\omega(\Spc)$; here the right-hand side denotes the enriched Hom-space as a left $\Cond^\omega(\Spc)$-module, see Appendix~\ref{app.yoneda}.
\end{thm}
\begin{proof}
First, we construct a map. For every affinoid algebra $A$ we have a map 
\[ \ZZ(A)\times\BGL(A) \To \Omega^\infty\K_{\geq 0}(A) \To \Omega^\infty\Kan_{\geq 0}(A) \]
which is functorial in $A$, hence inducing a map $\ZZ\times\BGL\to\Omega^\infty\Kan_{\geq 0}$ on the category of presheaves on affinoid rigid spaces. Thus we obtain an induced map $\Lmot(\ZZ\times\BGL) \to \Omega^\infty\Kan_{\geq 0}$ in $\RigH(S)$. 

Using the fibre sequence \eqref{fibre-sequence-1-connective-cover} we can separately investigate the induced maps $\Lmot\BGL\to\Omega^\infty\Kan_{\geq 1}$ on 1-connective covers and $\Lmot\ZZ\to\Kan_0$ in degree zero. On 1-connective covers, we compute for every affinoid algebra $A$ as follows.\footnote{Note that $\BGL$ arises from sheafifying the composition of the presheaf $\Rig_S\op \to \Spc,\,X\mapsto\BGL(\Ocal_X(X)),$ with the inclusion $\Spc\inj\Pro(\Spc)$.}
\begin{align*}
(\L_{\Arig}\BGL)(\Spa(A)) &\simeq \colim_{n\in\Delta\op}\, \BGL(A\skp{\Delta^n})\\
&\simeq \colim_{n\in\Delta\op}\, \Hom(\colim_\rho \Spa(A\skp{\Delta^n}_\rho),\BGL)\\
&\simeq \colim_{n\in\Delta\op}\, \lim_\rho \Hom(\Spa(A\skp{\Delta^n}_\rho),\BGL).\\
(\clubsuit) &\simeq \lim_{\rho\in\NN\op}\, \colim_{n\in\Delta\op}\, \Hom(\Spa(A\skp{\Delta^n}_\rho),\BGL)\\
&\simeq \lim_{\rho\in\NN\op}\, \BGL(A\skp{\Delta}_\rho) \\
&\simeq \KVan(A) \simeq \Omega^\infty\tau_{\geq 1}\Kan(A)
\end{align*}
where the equivalence $(\clubsuit)$ follows from \cite[Lemma 2.8]{kst-ii}.
In degree zero we have that $\Kan_0(A) \cong \Kcont_0(A) \cong \K_0(A) \cong \ZZ$ for local rings $A$ which shows the first claim.

It remains to show that
$$ 
\Omega^\infty\Kan_{\geq 0}(X) \simeq \Homline_{\Cond^\omega(\Spc)}(\Lmot X,\Lmot(\ZZ\times\BGL)).
$$
For this let us note that by construction, we have that $\Lmot$, as a localisation from\linebreak $\Fun(\SmRig_{S}\op,\Cond^{\omega}(\Spc))\to \RigH(S,\Cond^{\omega}(\Spc))$, is equivalently given by base\linebreak change of the motivic localisation functor $\PSh(\SmRig_S)\to \RigH(S)$. In particular, it is compatible with the $\Cond^{\omega}(\Spc
)$-module structure. Thus, we have 
$$
\Homline_{\Cond^\omega(\Spc)}(\Lmot X,\Lmot(\ZZ\times\BGL))\simeq \Homline_{\Cond^\omega(\Spc)}(X,\Lmot(\ZZ\times\BGL)),
$$
where we view $X$ as a functor in $\Fun(\SmRig_S,\Cond^\omega(\Spc))$ via the composition of the Yoneda functor and the constant functor $\Spc\to \Cond^\omega(\Spc)$.
Our results above together with the enriched Yoneda lemma finish the proof  (cf. Lemma \ref{lem-yoneda}).
\end{proof}

\begin{defi}
For $n\geq d\geq 1\in\ZZ$ denote by $\textup{Grass}_{d,n}$ be the Grassmannian $R$-scheme representing submodules $\Ecal\subseteq\Ocal^n$ such that the quotient sheaf $\Ocal^n/\Ecal$ is locally free of rank $n-d$ \cite[(8.4)]{GW1}. We write $\textup{Grass}_d := \colim_n\textup{Grass}_{d,n}$ and $\textup{Grass} := \colim_d\textup{Grass}_d$. Denote by $\textup{Grass}_{d,n}^\an$, $\textup{Grass}_d^\an$, and $\textup{Grass}^\an$ the respective analytifications.

Note that $\textup{Grass}_{d,n}^\an$ coincides with the representable functor $\textup{Grass}_{d-n}^\rig$ representing $d$-dimensional subspaces of $\kappa(x)^n$ on points $\Spa(\kappa(x))\to S$.
\end{defi}


\begin{prop}[Arndt, Sigloch {\cite[Prop.~5.34]{sigloch-thesis}}] \label{Prop:BGL=Grass}
For every $n\geq 1$, there is an equivalence $\BGL_n\simeq\textup{Grass}_n^\rig$ in $\RigH(S)$. Consequently, there is an equivalence $\BGL\simeq\textup{Grass}^\an$.
\end{prop}


\begin{cor}
\label{cor-algebraic-analytic-K}
We assume $(\dagger)_R$. Then the image of connective algebraic K-theory under the analytification functor $\alpha \colon \textup{H}(\Spec(R),\Spc)\to\RigH(S,\Cond^\omega(\Spc))$ identifies with connective analytic K-theory.
\end{cor}
\begin{proof}
In $\textup{H}(\Spec(R),\Spc)$, algebraic K-theory coincides with $\Lmot(\ZZ\times\BGL)$ \cite[Thm.~3.13]{MV1}. In $\RigH(S,\Cond^\omega(\Spc))$, connective analytic K-theory coincides with\linebreak  $\Lmot(\ZZ\times\BGL)$ (cf. Theorem~\ref{thm:representability}). The identification $\alpha(\ZZ)\simeq\ZZ$ is clear. The identification $\alpha(\BGL)\simeq\BGL$ follows since $\BGL\simeq\textup{Grass}$ on both sides by \cite[Prop.~3.7]{MV1} and Proposition~\ref{Prop:BGL=Grass}.
\end{proof}




%%%%%%%%%%%%%%%%%%%%%%%%%%%%%%%%%%%%%%%%%%%%%
\subsection{Analytic Bass delooping}
\label{sec:analytic-bass-construction}
%%%%%%%%%%%%%%%%%%%%%%%%%%%%%%%%%%%%%%%%%%%%%
Nonconnective algebraic K-theory is obtained from connective algebraic K-theory by Bass delooping within the category of spectra \cite[\S 6]{tt90}. The construction of the delooping uses the cover of $\PP^1$ by two copies of $\Arig$ whose intersection is $\Gm$. Such a delooping construction can also be done globally within the category of presheaves of spectra \cite[\S 2]{cisinski-descente}. For analytic K-theory of rigid spaces, an analytic analogue of the Bass construction is performed by Kerz-Saito-Tamme within the category of pro-spectra and using the cover of $\PP^1$ by two copies of $\Brig$ whose intersection we denote by $\HH_m$. Putting 1-connective analytic K-theory into this machine, one obtains, by definition, nonconnective analytic K-theory, see subsection~\ref{subsec.continuous-analytic-k-theory}.
In this subsection, we examine such an analytic Bass delooping within the category $\PSh(\Rig_S,\Pro(\Sp^+))$ using the alternative cover of $\PP^1$ by two copies of $\Arig$ whose intersection is $\Gm$, the analytification of $\Spec(R[t,t^{-1}])$. It turns out that the $\Arig$-delooping is equivalent to the $\Brig$-delooping so that  analytic K-theory also satisfies the Bass Fundamental Theorem for $\Arig$ (Corollary~\ref{bass-fundamental-theorem-for-Kan--cor}). Our exposition follows closely the description of the $\Brig$-analytic Bass delooping by Kerz-Saito-Tamme \cite[\S 4.4]{kst-i}.

\vspace{6pt}\noindent\textbf{Notation.} We fix an embedding $\Arig \inj \PP^1$ and write $\PP^1\setminus\Arig=:\{\infty\}$ and $\AA^{-1}:=\PP^1\setminus\{O\}$ where $O\in\Arig$ denotes the origin. Thus $\Gm=\Arig\cap\AA^{-1}$ as subspaces of $\PP^1$. Analogously, we write $\Brig:=\Spa(R\skp{t},R^+\skp{t}) \inj \PP^1$ and $\BB^{-1}:=\Spa(R\skp{t^{-1}},R^+\skp{t^{-1}})$, so that $\Brig\cap\BB^{-1}=\Spa(R\skp{t,t^{-1}},R^+\skp{t,t^{-1}})=:\HH_\textup{m}$.
Finally, let $\Dcal$ be a stable presentable category; thus it carries an essentially unique structure of a module over the category of spectra \cite[4.8.2.18]{HA}.

\begin{defi} \label{bass-delooping--defi}
Let $E\in\PSh(\Rig_S,\Dcal)$ be a presheaf. We define a presheaf $\Gamma_{\Arig}E$  via
\[ \Gamma_{\Arig}E(X) := E(X\times\Arig) \sqcup_{E(X)} E(X\times\AA^{-1}). \]
The commutativity of the square
\carre{\Gm}{\Arig}{\AA^{-1}}{\PP^1}
yields a map $\Gamma_{\Arig}E(X) \to E(X\times \Gm)$ so that we can define a presheaf $\Lambda_{\Arig}E$ via
\[ \Lambda_{\Arig}E(X) := \textup{fib}\bigl(\Gamma_{\Arig}E(X) \to E(X\times\Gm)\bigr). \]

Replacing $\Arig$, $\AA^{-1}$, and $\Gm$ with $\Brig$, $\BB^{-1}$, and $\HH_\textup{m}$, respectively, we get the presheaves $\Gamma_{\Brig}E$ and $\Lambda_{\Brig}E$ as defined in the affinoid setting by Kerz-Saito-Tamme \cite[\S 4.4]{kst-i}. 
\end{defi}

\begin{lem} \label{bass-constructions-agree--lem}
Given a presheaf $E\in\PSh(\Rig_S,\Dcal)$, there are equivalences
\[ \Lambda_{\Arig}E(X) \simeq \Sigma\textup{fib}\bigl(E(X\times\PP^1) \to E(X)\bigr) \simeq \Lambda_{\Brig}E(X) \]
which are functorial in $X\in\Rig_S$.
\end{lem}
\begin{proof}
We have a map of cartesian squares
$$ \begin{tikzcd}
    E(X) \arrow[r] \arrow[dd] 
    & E(X\times\Arig) \arrow[dd]
    &&
    E(X\times\PP^1) \arrow[r] \arrow[dd] 
    & E(X\times\Arig) \arrow[dd]
    \\ && \overset{\alpha}{\longrightarrow} \\
    E(X\times\AA^{-1}) \arrow[r] 
    &\Gamma_{\Arig}E(X) 
    &&
    E(X\times\AA^{-1}) \arrow[r] 
    & E(X\times\Gm).
\end{tikzcd} $$
Denoting $\Phi E(X) := \textup{fib}\bigl(E(X\times\PP^1) \to E(X)\bigr)$, the fibre of the map $\alpha$ is the cartesian square
\carre{\Phi E(X)}{\ast}{\ast}{\Lambda_{\Arig}E(X).}
This shows the first asserted equivalence. The same argument with $\Gamma_{\Brig}E$ and $\Lambda_{\Brig}E$ yields the second asserted equivalence.
\end{proof}

\begin{defi} \label{bott-element--defi}
Let $E\in\PSh(\Rig_S,\Dcal)$ carrying a $\K$-module structure. We fix a map $t\colon\SS\to\Omega\K(\ZZ[t,t^{-1}])$ representing the class $t\in\K_1(\ZZ[t,t^{-1}])$. Then we get an induced map
\[ \lambda \colon E(X) \To[-\cup t] \Omega E(X\times\Gm) \To[\partial] \Lambda_{\Arig}E(X) \]
where the first map comes from the cup product on K-theory and the second map is from the shifted fibre sequence defining $\Lambda_{\Arig}E$ above.
\end{defi}

\begin{defi} \label{analytic-bass-construction--defi}
Let $E\in\PSh(\Rig_S,\Dcal)$ carrying a $\K$-module structure. We write $\Lambda E := \Lambda_{\Arig}E \simeq\Lambda_{\Brig}E$ and define the analytic \emph{analytic Bass construction of} $E$  as the colimit
\[ E^\textup{B} := \colim\bigl( E \To[\lambda] \Lambda E \to[\lambda] \Lambda^2 E \To[\lambda] \ldots \bigr) \]
within the category $\PSh(\Rig_S,\Dcal)$.
\end{defi}


As a formal consequence from our construction we get the following.

\begin{prop} \label{bass-fundamental-theorem-abstractly--prop}
Let $\Dcal\in\{\Sp,\Pro(\Sp),\Cond^\omega(\Sp)\}$.
Let $E\in\PSh(\Rig_S,\Dcal)$ carrying a $\K$-module structure and assume that the map $\lambda \colon E\to \Lambda E$ is an equivalence. Then for every $n\in\ZZ$ and every $X\in\Rig_S$ we have an exact sequence
\[ 0 \To E_n(X) \To E_n(X\times\Arig) \oplus E_n(X\times\AA^{-1}) \To[\pm] E_{n}(X\times\Gm) \To[\partial] E_{n-1}(X) \To 0 \]
in the category $\Dcal^\heartsuit$. Furthermore, the map $\partial$ is split by the map
\[
-\cup t \,\colon\, E_{n-1}(X) \To E_n(X\times\Gm) 
\]
induced by the cup product with $t$ from Definition~\ref{bott-element--defi}.
\end{prop}
\begin{proof}
The same argument as for the $\Brig$-delooping applies, see~\cite[Prop.~4.13]{kst-i}.
\end{proof}

\begin{cor} \label{bass-fundamental-theorem-for-Kan--cor}
For every $n\in\ZZ$ and every $X\in\Rig_S$ there is an exact sequence
\[ 0 \To \Kan_n(X) \To \Kan_n(X\times\Arig) \oplus \Kan_n(X\times\AA^{-1}) \To[\pm] \Kan_{n}(X\times\Gm) \To[\partial] \Kan_{n-1}(X) \to 0 \]
of pro-abelian groups where the map $\partial$ has a split.
\end{cor}
\begin{proof}
This follows from Proposition~\ref{bass-fundamental-theorem-abstractly--prop} since the map $\lambda \colon \Kan \to \Lambda\Kan$ is an equivalence \cite[Cor.~2.6]{kst-ii}.
\end{proof}


\begin{cor}
\label{cor-K-P1-spectrum}
  There exists a $\PP^1$-action on $\ZZ\times\BGL$ such that the induced $\PP^1$-spectrum
    $$
    \textup{KGL}^{\an}\coloneqq (\ZZ\times\BGL,\ZZ\times\BGL,\dots)\in \RigSH(K)
    $$
    is a commutative algebra object in $\RigSH(K)$. Furthermore, under assumption $(\dagger)$ the $\PP^1$-spectrum $\textup{KGL}^{\an}$ represents (non-connective) analytic $K$-theory.
\end{cor}
\begin{proof}
    Classically, it is known that $\ZZ\times \BGL\in \textup{H}(K)$, admits a $\PP^1$-action \cite[\S 6.2]{voevodsky-icm-1998}, \cite{PPR-KGL}, \cite[\S 3.2]{MotCoh}.
    This action yields a $\PP^1$-spectrum $\textup{KGL}\coloneqq(\ZZ\times \BGL,\ZZ\times \BGL,\dots)$ representing homotopy invariant K-theory. By our discussion in Section \ref{sec-RigSH-as-spectra}, we see that the analytification of $\textup{KGL}$ yields a $\PP^1$-spectrum $\textup{KGL}^{\an}$ inside\linebreak $\RigSH(K,\Cond^{\omega}(\Sp^+))$. Furthermore, by Lemma \ref{bass-constructions-agree--lem} and the construction of homotopy invariant K-theory ring spectrum via the Bass-construction \cite{k-book}, we see that $\textup{KGL}^{\an}$ represents (non-connective) analytic K-theory under assumption $(\dagger)$ using Theorem \ref{thm:representability}.
\end{proof}


%%%%%%%%%%%%%%%%%%%%%%%%%%%%%%%%%%%%%%%%%%%%%
\subsection{Continuous Homotopy K-theory}
%%%%%%%%%%%%%%%%%%%%%%%%%%%%%%%%%%%%%%%%%%%%%

\vspace{6pt}\noindent Under certain finiteness conditions, analytic K-theory is the $\Arig$-localisation of continuous K-theory.

\begin{thm}[$\Arig$-localisation]
\label{A1-localisation--thm}
Let $A$ be a Tate ring that admits a noetherian and finite dimensional ring of definition. Then there exists a canonical weak equivalence
\[ (\L_{\Arig}\Kcont)(A) \simeq \Kan(A). \]
\end{thm}
\begin{proof}
We have the following chain of equivalences:
\begin{align*}
(\L_{\Arig}\Kcont)(A) 
&\simeq \colim_{n\in\Delta\op} \Kcont(\Delta_A^{\an,n})
&(\textup{Lemma}~\ref{homotopy-invariantification--lem})  \\
&\simeq \colim_{n\in\Delta\op} \prolim_j \Kcont(A\skp{\Delta_{\pi^j}})
&(\textup{sheaf condition}) \\
&\simeq \prolim_j \colim_{n\in\Delta\op} \Kcont(A\skp{\Delta_{\pi^j}})
& (\clubsuit) \\
&\simeq \prolim_j \Kcont(A\skp{\Delta_{\pi^j}})
& (\textup{definition}) \\
&\simeq \Kan(A) 
& \cite[\textup{Cor.}~2.9]{kst-ii}
\end{align*}
For the equivalence $(\clubsuit)$ we use that the geometric realisation commutes with the pro-limit \cite[\textup{Lem.}~2.8]{kst-ii} since the spectra in question are uniformly bounded below by \cite[Thm.~8.12]{dahli-thesis-paper}.
\end{proof}

Continuous K-theory seems to be the canonical K-theory for rigid analytic spaces in the view of the Efimov Continuity Theorem \cite[Satz~5.8]{andreychev-thesis}. Thus, in analogy to Weibel's homotopy K-theory, its $\Arig$-localisation is a natural cohomology theory to study, whence the following.

\begin{defi} \label{defi-continuous-homotopy-k-theory}
We define \emph{continuous homotopy K-theory} to be the $\Arig$-localisation of continuous K-theory, i.e. $\KHcont := \L_{\Arig}\Kcont$.
\end{defi}

\begin{lem} \label{lem--KHcont-is-Nisnevich-sheaf}
The functor $\KHcont \colon \Rig\to\Pro(\Sp^+)$ is a Nisnevich sheaf. In particular, $\KHcont \simeq \Lmot\Kcont$.
\end{lem}
\begin{proof}
Let $U\to X$ be a Nisnevich cover so that $X\simeq\colim_{i\in\Delta\op}U_i$ where $U_i$ is the $(i+1)$-fold fibre product of $U$ over $X$. We have the following equivalences
\begin{align*}
\KHcont(X)
&\simeq \colim_{n\in\Delta\op} \Kcont(X\times \Delta^{\an,n}) 
& (\textup{Definition}~\ref{defi-continuous-homotopy-k-theory}) \\
&\simeq \colim_{n\in\Delta\op} \prolim_i \Kcont(U_i\times\Delta^{\an,n}) 
& (\textup{sheaf condition}) \\
&\simeq \prolim_i \colim_{n\in\Delta\op} \Kcont(U_i\times\Delta^{\an,n}) 
& \cite[\textup{Cor.}~2.9]{kst-ii} \\
& \simeq \prolim_i \KHcont(X)
& (\textup{Definition}~\ref{defi-continuous-homotopy-k-theory})
\end{align*}
which shows the claim.
\end{proof}




%%%%%%%%%%%%%%%%%%%%%%%%%%%%%%%%%%%%%%%%%%%%%
\begin{appendix}
%%%%%%%%%%%%%%%%%%%%%%%%%%%%%%%%%%%%%%%%%%%%%

%%%%%%%%%%%%%%%%%%%%%%%%%%%%%%%%%%%%%%%%%%%%%%%%%%%%%%
\section{Condensed objects}
%%%%%%%%%%%%%%%%%%%%%%%%%%%%%%%%%%%%%%%%%%%%%%%%%%%%%%

In the main body of this article, we deal with analytic K-theory as defined and studied by Kerz-Saito-Tamme \cite{kst-i,kst-ii} which is a functor
\[ \Kan \colon \Rig \To \Pro(\Sp) \]
with values in the category of pro-spectra. Since pro-categories are not that well-behaved, e.g.\ computing colimits in them is not easy, we eventually will postcompose $\Kan$ with the canonical functor $\Pro(\Sp)\to\Cond(\Sp)$ to condensed spectra.
%First, we start with a brief recollection about pro-objects.

\vspace{6pt}\noindent\textbf{Notation.} For this entire section, we fix an uncountable strong limit cardinal $\kappa$.
%For a category $\Dcal$ we consider its associated category $\Cond_\kappa(\Dcal)$ of $\kappa$-condensed objects as well as the category $\Cond^\omega(\Dcal)$ of light condensed objects. 





%%%%%%%%%%%%%%%%%%%%%%%%%%
\subsection{Condensed objects in a category}
%%%%%%%%%%%%%%%%%%%%%%%%%%

Given a category $\Ccal$, we work with the categories $\Cond^\omega(\Ccal)$ of light condensed objects in $\Ccal$ and $\Cond_\kappa(\Ccal)$ of $\kappa$-condensed objects in $\Ccal$, as opposed to the large category $\Cond(\Ccal) := \colim_\lambda\Cond_\lambda(\Ccal)$ of condensed objects (where $\lambda$ runs over all uncountable strong limit cardinals). One advantage is the following: If the category $\Ccal$ is presentable, then $\Cond^\omega(\Ccal)$ and $\Cond_\kappa(\Ccal)$ are also presentable (Lemma~\ref{condensed-objects-presentable--lem}), whereas the category $\Cond(\Ccal)$ is not. Hence we can apply Adjoint Functor Theorems, for instance. The material about light condensed objects is extracted from the lectures by Clausen and Scholze on Analytic Stacks.\footnote{See \url{https://www.youtube.com/playlist?list=PLx5f8IelFRgGmu6gmL-Kf_Rl_6Mm7juZO}.}

\begin{defi}[Condensed objects]
\begin{enumerate}
\item A profinite set is called \emph{light} if it is isomorphic to a countable limit of finite sets. Denote by $\ProFin^\omega$ the full subcategory of $\ProFin$ that is spanned by light profinite sets. 
We equip the category $\ProFin^\omega$ with the topology generated by finite families of jointly surjective maps.
\item Analogously, let $\ProFin_{<\kappa}$ be the full subcategory of $\ProFin$ that is spanned by $\kappa$-small sets, equipped with the topology generated by finite families of jointly surjective maps.
\end{enumerate}
Let $\Ccal$ be a category that admits finite limits.
\begin{enumerate} \setcounter{enumi}{2}
\item A \emph{light condensed object in} $\Ccal$ is a sheaf $X$ on $\ProFin^\omega$ with values in $\Ccal$.
We denote $\Cond^\omega(\Ccal) := \Sh(\ProFin^\omega,\Ccal)$ the category of light condensed objects in $\Ccal$.
\item A $\kappa$\emph{-small condensed object in} $\Ccal$ is a sheaf on $\ProFin_{<\kappa}$ with values in $\Ccal$.
We denote $\Cond_\kappa(\Ccal) := \Sh(\ProFin_{<\kappa} ,\Ccal)$ the category of $\kappa$-small condensed objects in $\Ccal$.
\end{enumerate}
\end{defi}

\begin{rem}
Explicitly, a presheaf $X\colon (\ProFin^\omega)\op\to\Ccal$ is a sheaf if and only if it satisfies
\begin{enumerate}
    \item $X(\varnothing)\simeq *$,
    \item $X(S\sqcup S') \to[\simeq] X(S)\times X(S')$, and
    \item $X(S) \to[\simeq] \lim_\Delta X(\Check{C}(T\to S))$ 
\end{enumerate}
for all $S,S'$ and every surjective map $T\surj S$ in $\ProFin^\omega$.
\end{rem}



\begin{lemma} \label{condensed-objects-presentable--lem}
Let $\Ccal$ be a presentable category. Then the categories $\Cond^\omega(\Ccal)$ and $\Cond_\kappa(\Ccal)$ are presentable.
\end{lemma}
\begin{proof}
Being a category of sheaves on the small site $\ProFin^\omega$, the category $\Cond^\omega(\Spc)$ is an $\infty$-topos \cite[6.2.2.7]{HTT}. We have an equivalence $\Cond^\omega(\Ccal) \simeq \Sh(\Cond^\omega(\Spc),\Ccal)$ \cite[1.3.1.7]{SAG} and the latter category is presentable as $\Ccal$ is presentable \cite[1.3.1.6]{SAG}. 
Same argument for $\Cond_\kappa(\Ccal)$.
\end{proof}

\begin{lemma} \label{condensed-objects-stable--lem}
Let $\Ccal$ be a stable category. Then the categories $\Cond^\omega(\Ccal)$ and $\Cond_{<\kappa}(\Ccal)$ are stable.
\end{lemma}
\begin{proof}
As $\Ccal$ is stable, the functor category $\Fun(\ProFin^{\omega,\opp},\Ccal)$ is stable \cite[1.1.3.1]{HA}.  Since the sheafification functor $\Fun(\ProFin^{\omega,\opp},\Ccal)\to\Cond^\omega(\Ccal)$ is left-exact, the inclusion functor $\Cond^\omega(\Ccal) \inj \Fun(\ProFin^{\omega,\opp},\Ccal)$ creates finite limits. Hence $\Cond^\omega(\Ccal)$ admits finite limits and the suspension functor is an equivalence, thus the category is stable \cite[1.4.2.27]{HA}.
Same argument for $\Cond_\kappa(\Ccal)$.
\end{proof}



\begin{lemma} \label{condensed-objects-symmetric-monoidal}
Let $\Ccal$ be a category that admits small colimits and finite limits. Then the categories $\Cond^\omega(\Ccal)$ and $\Cond_\kappa(\Ccal)$ both admit finite limits. In particular, they admit the structure of a cartesian symmetric monoidal category. 
\end{lemma}
\begin{proof}
By assumption on $\Ccal$, the functor category $\Fun(\ProFin^{\omega,\opp},\Ccal)$ admits small colimits and finite limits \cite[5.1.2.3]{HTT} and this property passes to its full subcategory $\Cond^\omega(\Ccal)$ of sheaves.
Now the statement about the cartesian symmetric monoidal category from the existence of finite products \cite[2.4.1.5.(5)]{HA}.
Same argument for $\Cond_\kappa(\Ccal)$.
\end{proof}


\begin{remark}[Extremally disconnected sets]
\label{extremally-disconnected--rem}
A profinite set is called an \emph{extremally disconnected set} if it is a projective object in the category $\ProFin$. Denote by $\EDS_{<\kappa}$ the full subcategory of $\ProFin_{<\kappa}$ spanned by extremally disconnected sets. By definition of being projective, every cover of profinite sets $S'\to S$ with $S$ extremally disconnected has a split. In particular, the identity map $\id_S$ is cofinal among all covers of $S$. We shall use this property in the proof of Lemma~\ref{pro-vanishing-and-Mittag-Leffler--lemma} below.

As a matter of fact, the restriction along the inclusion functor $\EDS_{<\kappa} \inj \ProFin_{<\kappa}$ induces an equivalence from the category $\Cond_\kappa(\Ccal)$ onto the full subcategory\linebreak $\Fun^\times(\EDS_{<\kappa}\op,\Ccal)$ of $\Fun(\EDS_{<\kappa}\op,\Ccal)$ which is spanned by those objects that preserve finite products, i.e.\ sending finite disjoint unions in $\EDS_{<\kappa}$ to products in $\Ccal$.
\end{remark}

\begin{remark}[light vs. $\kappa$-small] 
\label{comparison-light-kappa-condensed}
Let $\Ccal$ be a category that admits limits and colimits. The restriction along the inclusion functor $j\colon \ProFin^\omega \inj \ProFin_{<\kappa}$ induces an adjunction

\begin{center}
\begin{tikzcd}
 \Fun(\ProFin^{\omega,\opp},\Ccal) \arrow[r, bend left, "j_!"] \arrow[r, bend right, "j_*", swap] &  \Fun(\ProFin_{<\kappa}\op,\Ccal) \arrow[l, "j^*", swap]
\end{tikzcd}
\end{center}
where $j_!$ is a left Kan extension and $R:=j_*$ is a right Kan extension.
 %    \[
	%    j^* \colon \Fun(\ProFin_{<\kappa}\op,\Ccal) \rightleftarrows \Fun(\ProFin^{\omega,\opp},\Ccal) \colon R
	% \]
%where $R$ is a right Kan extension.
Hence the functor $j^*$ preserves limits, so that it preserves sheaves and restricts to a functor $j^* \colon \Cond_\kappa(\Ccal)\to\Cond^\omega(\Ccal)$. Now the localisation functor $\Fun(\ProFin_{<\kappa}\op,\Ccal) \to \Cond_\kappa(\Ccal) \simeq \Fun^{\times}(\EDS\op_{<\kappa},\Ccal)$ factors over the restriction functor $\Fun(\ProFin_{<\kappa}\op,\Ccal) \to\Fun(\EDS_{<\kappa}\op,\Ccal)$. Given a light condensed object $X\in\Cond^\omega(\Ccal)$ and $S_1,S_2\in\EDS_{<\kappa}$, we have
\begin{align*}
    RX(S_1\sqcup S_2) 
    &\simeq \lim_{\ProFin^\omega\ni W\to S_1\sqcup S_2} X(W) \\
    &\simeq \lim_{\ProFin^\omega\ni W_1\to S_1} X(W_1) \times \lim_{\ProFin^\omega\ni W_2\to S_2} X(W_2) \\
    &\simeq RX(S_1) \times RX(S_2).
\end{align*} 
Hence the adjunctions above induce adjunctions 
%\footnote{Note that $\Cond_{\kappa}(\Ccal)$ is equivalent to the category of finite product preserving functors $F\colon \EDS_{<\kappa}\op\to \Ccal$ and the inclusion $\EDS_{<\kappa}\hookrightarrow \ProFin_{<\kappa}$ preserves finite coproducts.}
\begin{center}
\begin{tikzcd}
 \Cond^{\omega}(\Ccal) \arrow[r, bend left, "\L_\kappa j_!"] \arrow[r, bend right, "j_*", swap] &  \Cond_{\kappa}(\Ccal) \arrow[l, "j^*", swap]
\end{tikzcd}
\end{center}
where $j^*$ and $j_*$ are just the restricted functors.
% \[
% 	   j^* \colon \Cond_{\kappa}(\Ccal) \rightleftarrows \Cond^{\omega}(\Ccal) \colon R.
% 	\]
Via the explicit formula for Kan extensions we see that the composition $j^*\circ R$ is equivalent to the identity functor, hence $R$ is fully faithful.
Since colimits commute with finite products we get analogously as above that
\begin{align*}
    j_!X(S_1\sqcup S_2) 
    &\simeq \colim_{S_1\sqcup S_2 \to W\in\ProFin^\omega} X(W) \\
    &\simeq \colim_{S_1 \to W_1\in\ProFin^\omega} X(W_1) \times \colim_{S_2 \to W_2\in\ProFin^\omega} X(W_2) \\
    &\simeq j_!X(S_1) \times j_!X(S_2).
\end{align*} 
\end{remark} 




%%%%%%%%%%%%%%%
\subsection{From pro-objects to condensed objects}
%%%%%%%%%%%%%%%



\textbf{Notation.} In this subsection, let $\Ccal$  be a category that is accessible, admits colimits and admits finite limits.
Denote by $\Pro(\Ccal)$ the associated category of pro-objects $\Ccal$, cf.\ the beginning of section~\ref{pro-categories--subsubsection}.
We have adjunctions
	\[
	\textup{const} \colon \Ccal \rightleftarrows \Pro(\Ccal) \colon \textup{lim}\quad\quad\quad \textup{const}\colon\Ccal \rightleftarrows \Cond^{\omega}(\Ccal) \colon\textup{can}
	\]
where the left-adjoints are the respective constant functors which are fully faithful.

\begin{lemma} \label{comparison-functor--lemma}
Assume additionally that $\Ccal$ admits small cofiltered limits\footnote{In particular, $\Ccal$ admits all limits since it has pullbacks and products (which are cofiltered limits of finite products) \cite[Prop. 4.4.2.6]{HTT}.}. Then there is a canonical comparison functor 
\[ \gamma_\kappa \colon \Pro(\Ccal) \to \Cond_{\kappa}(\Ccal) \]
commuting with small limits and fitting into a commutative diagram
\[ \begin{xy} \xymatrix{
	&\Ccal \ar[dl]_{\mathrm{const}} \ar[d]^{\mathrm{const}} \ar[dr]^{\mathrm{const}} \\
	\Pro(\Ccal) \ar[dr]_{\mathrm{lim}} \ar[r]^-{\gamma_\kappa} & \Cond_{\kappa} \ar[d]^{\mathrm{can}}(\Ccal) \ar[r]^{j^*} & \Cond^\omega(\Ccal) \ar[dl]^{\mathrm{can}}
	 \\ & \Ccal
} \end{xy} \]
Explicitly, the functor $\gamma^{\kappa}$ sends the pro-object $\prolim_{i\in I}X_i$ to the limit $\lim_{i\in I}\underline{X_i}$.
%The analogous result for $\Cond_\kappa(\Ccal)$ is true as well.
\end{lemma}
\begin{proof}
The commutativity of the lower right triangle follows since can is given by evaluating at a single point.
For showing the commutativity of the upper right triangle consider the following square of left adjoint functors
\[ \begin{tikzcd}
    \Fun(\ProFin_{<\kappa}\op,\Ccal) \arrow[d,"\L_\kappa",swap]\arrow[r,"j^*"]
    &  \Fun(\ProFin^\omega,\Ccal) \arrow[d,"\L^\omega"]\\
    \Cond_\kappa(\Ccal) \arrow[r,"j^*"] & \Cond^{\omega,\opp}(\Ccal).
\end{tikzcd} \]
This square commutes because the associated square of right adjoint functors commutes.
Since $\Ccal$ admits small filtered limits, the category $\Cond_{\kappa}(\Ccal)$ admits small filtered limits as well. Hence restriction along the constant functor $\Ccal \to \Pro(\Ccal)$ induces an equivalence 
\[ \Fun^\mathrm{lim}(\Pro(\Ccal),\Cond_\kappa(\Ccal)) \To[\simeq] \Fun(\Ccal,\Cond_\kappa(\Ccal)) \]
where the source denotes the full subcategory of all functors spanned by those which preserve small filtered limits \cite[A.8.1.6]{SAG}. Thus the constant functor $\Ccal \to \Cond_\kappa(\Ccal)$ yields the functor $\gamma_\kappa$ and the upper left commutative triangle. 
Analogously, we get get a functor $\gamma^\omega \colon \Pro(\Ccal) \to \Cond^\omega(\Ccal)$ restricting to the constant functor $\Ccal \to \Cond^\omega(\Ccal)$. Since both $\gamma^\omega$ and $j^*\circ\gamma_\kappa$ restrict along to the constant functor, we see that they are equivalent.
%Analogously, we have an equivalence
%\[ \Fun'(\Pro(\Ccal),\Ccal)) \To[\simeq] \Fun(\Ccal,\Ccal) \]
%\hl{yielding the lower commutative triangle.} 

The commutativity of the lower left triangle follows if $\gamma := \gamma_\kappa$ preserves all limits, as both compositions $\lim\circ \mathrm{ const}$ are equivalent to the identity functor.


We are left to show that $\gamma$ preserves finite pullbacks, as then $\gamma$ preserves all limits by \cite[Prop. 4.4.2.7]{HTT}. For this, we remark that finite limits in $\Pro(\Ccal)$ can be computed pointwise in the sense of \cite[\S 2.1]{kst-i}. Thus, any diagram 
$$
\begin{tikzcd}
    & X\arrow[d,""]\\
    Y\arrow[r,""]& Z
\end{tikzcd}
$$
in $\Pro(\Ccal)$ can be lifted to an element $F\in \Pro(\Fun(K,\Ccal))$, where $K$ is the set representing the diagram above, such that its image under 
$$
    \Pro(\lim)\colon \Pro(\Fun(K,\Ccal))\To\Pro(\Ccal)
$$
is equivalent to $X\times_Z Y$. We also have a functor
$$
    \Cond(\lim)\colon \Cond_\kappa(\Fun(K,\Ccal))\To\Cond_\kappa(\Ccal)
$$
that is induced by the following construction. We have a limit preserving functor 
$$ \gamma^K\colon \Cond_\kappa(\Fun(K,\Ccal))\inj\Fun(\ProFin_{<\kappa},\Fun(K,\Ccal))\simeq \Fun(K,\Fun(\ProFin_{<\kappa},\Ccal)).
$$
Then $\Cond(\lim)$ is the composition of the above functor with the limit functor and sheafification. Now we obtain an essentially commutative square
\begin{equation}\tag{$\clubsuit$}   
    \begin{tikzcd}
        \Pro(\Fun(K,\Ccal)) \arrow[d,"\Pro(\lim)",swap]\arrow[r,"\gamma^{K}"]& \Cond_\kappa(\Fun(K,\Ccal))\arrow[d,"\Cond(\lim)"]\\
        \Pro(\Ccal)\arrow[r,"\gamma"]&\Cond_\kappa(\Ccal).
    \end{tikzcd}
\end{equation}
Note that limits and colimits in functor categories can be computed pointwise, thus we have an equivalence $\Cond_\kappa(\Fun(K,\Ccal))\simeq \Fun(K,\Cond_\kappa(\Ccal))$. Under this equivalence $\Cond(\lim)$ corresponds to the limit functor $\Fun(K,\Cond_{\kappa}(\Ccal))\rightarrow \Cond_\kappa(\Ccal)$. Therefore, commutativity of $(\clubsuit)$ shows that $\gamma$ preserves $K$-indexed limits, i.e. pullbacks. 
%Same argument for $\Cond_\kappa(\Ccal)$.
\end{proof}


\begin{remark} \label{continuous-comparison-lemma}
An informal argument for the comparison functor $\gamma$ to commute with cofiltered limits goes as follows:
Let $A \to \Pro(\Sp), \alpha \mapsto X_\alpha$ be a cofiltered diagram of pro-spectra $X_\alpha = \prolim_{i\in I_\alpha}X_{\alpha, i}$. We may assume that $A$ is cofinite and that the diagram $A\to\Pro(\Sp)$ is given in level-representation with cofiltered index set $I$ ($=I_\alpha$ for all $\alpha\in A$) so that $A\times I$ is cofiltered again and $\prolim_{A\times I}X_{\alpha,i} \simeq \lim_AX_\alpha$ is the limit \cite[\S 4]{isaksen-calculating}. Then
	\begin{align*}
	\gamma( \lim_{A} X_\alpha ) 
	\simeq \gamma( \prolim_{A\times I} X_{\alpha, i} ) 
	= \lim_{A\times I} \underline{X_{\alpha, i}} 
	\simeq   \lim_{A} \, \lim_{I} \underline{X_{\alpha, i}} 
	=  \lim_{A} \gamma( X_\alpha ).
	\end{align*}
\end{remark}



\begin{cor}
The comparison functor $\gamma$ preserves cofibre sequences and pushouts.
\end{cor}
\begin{proof}
This follows since both $\Pro(\Sp)$ and $\Cond(\Sp)$ are stable categories.
\end{proof}

Next we want to show that the comparison functor is conservative. For this purpose, we start with some preparations.

\begin{lemma} \label{condensed-milnor-sequence--lemma}
Let $\bigl( (X_i)_i, (p_i\colon X_i\to X_{i-1})_i \bigr)$ be a tower of spectra and $n\in\ZZ$. Then there is an exact sequence
\[
0 \To {\lim_i}^1\pi_{n+1}(X_i) \To \pi_n(\lim_iX_i) \To \lim_i\pi_n(X_i) \To 0
\]
of condensed abelian groups, the so-called Milnor sequence.
\end{lemma}
\begin{proof}
This is a standard proof.\footnote{Cf.~\url{https://ncatlab.org/nlab/show/lim^1+and+Milnor+sequences}.}
The sequential limit $\lim_iX_i$ fits into a fibre sequence
\[
\lim_i X_i \To \prod_iX_i \To[\id-p_*] \prod_iX_i
\]
where $p_*$ is the map induced by the maps $(p_i\colon X_i\to X_{i-1})_i$. Consider the associated long exact sequence of homotopy groups
\[
\ldots \To \prod_i\pi_{n+1}(X_i) \To[\id-p_*] \prod_i\pi_{n+1}(X_i) \To[\alpha] \pi_n(\lim_iX_i) \To[\beta] \prod_i\pi_n(X_i) \To[\id-p_*] \prod_i\pi_n(X_i) \To \ldots
\]
We have that
\[
\im(\beta) = \ker\bigl( \prod_i\pi_n(X_i) \To[\id-p_*] \prod_i\pi_n(X_i) \bigr) = \lim_i\pi_n(X_i)
\]
and
\[
\ker(\alpha) = \coker\bigl( \prod_i\pi_{n+1}(X_i) \To[\id-p_*] \prod_i\pi_{n+1}(X_i) \bigr) = {\lim_i}^1\pi_{n+1}(X_i)
\]
which shows the claim.
\end{proof}



The following result and its proof was communicated to us by Peter Scholze and we do not claim originality. 
%Any possible faults lie in the responsibility of the authors of this article.

\begin{thm}[Clausen-Scholze] \label{comparison-conservative--thm}
The comparison functor 
\[ \gamma_\kappa := R\circ\gamma^\omega \,\colon\, \Pro^\omega(\Sp^+) \to \Cond_\kappa(\Sp) \]
is conservative.
In particular, the comparison functor $\gamma^\omega$ is conservative, too.

\end{thm}
\begin{proof}
Since the categories in question are stable (Lemma~\ref{Pro-omega-stable--lemma} and Lemma~\ref{condensed-objects-stable--lem}), it suffices to check that the functor $\gamma$ reflects zero objects. Let $\ldots\to X_2 \to X_1 \to X_0$ be a
system of spectra such that $\lim_i\underline{X_i} = 0$ as condensed spectra (where each
$X_i$ is considered as a discrete/constant condensed spectrum, and the
limit is taken in condensed spectra). 
We have to see that the pro-spectrum $\prolim_iX_i$ is zero; this we can show on homotopy pro-groups $\lim_i\pi_j(X_i)$ for all $j\in\ZZ$ \cite[2.8]{kst-i} and we may assume $j=0$ by shifting. The Milnor sequences (Lemma~\ref{condensed-milnor-sequence--lemma}) for $n\in\{-1,0\}$
\[
0 \To {\lim_i}^1\pi_{n+1}(X_i) \To \pi_n(\lim_iX_i) \To \lim_i\pi_n(X_i) \To 0
\]
in $\Cond(\Ab)$ yield that $\lim_i \pi_0 X_i = \lim_i^1 \pi_0 X_i = 0$ in this category. Thus, the claim follows from the Lemma~\ref{pro-vanishing-and-Mittag-Leffler--lemma} below.
\end{proof}


\begin{remark}
The comparison functor $\gamma\colon\Pro^\omega(\Ab)\to\Cond_\kappa(\Ab)$ is not conservative. Consider the tower $(p^i\ZZ)_i$ whose structure maps are the inclusions (for a prime number $p$). Then $\prolim_ip^i\ZZ$ is not zero in $\Pro(\Ab)$ as for every $j\geq i$ the map $p^j\ZZ \to p^i\ZZ$ is not zero. But $\gamma(\prolim_ip^i\ZZ) = \lim_i\underline{p^i\ZZ} = \underline{\lim_ip^i\ZZ} =0$ since we can check this after passing to $\Top\Ab$ (Remark~\ref{lim-and-lim^1-in-Ab-and-CondAb--remark}) and since $\lim_ip^i\ZZ=0$ in $\Top\Ab$ as its underlying abelian group is zero.
\end{remark}

\begin{remark} \label{lim-and-lim^1-in-Ab-and-CondAb--remark}
The faithful functor $\Top\Ab \to \Cond_\kappa(\Ab), M\mapsto \underline{M} := C(-,M)$, commutes with limits and hence admits a left adjoint $\L \colon \Cond_\kappa(\Ab)\to\Top\Ab$. Thus for a tower $(M_i)_i$ of topological abelian groups with limit $M:=\lim_iM_i$ in $\Top\Ab$, the canonical morphism $\underline{M} \to \lim_i\underline{M_i}$ of associated condensed abelian groups is an isomorphism.
In $\Cond_\kappa(\Ab)$ we have an exact sequence
\[ \tag{$\clubsuit$}
0 \To \lim_i\underline{M_i} \To \prod_i\underline{M_i} \To[\delta] \prod_i\underline{M_i} \To {\lim_i}^1\underline{M_i} \To 0
\]
where $\delta$ is the difference of the identity and the structure map of the tower.
As the canonical morphism $\underline{\prod_iM_i}\To\prod_i\underline{M_i}$ is an isomorphism and since for every $N\in\Top\Ab$ the counit $\L(\underline{N})\to N$ is an isomorphism, the sequence $\L(\clubsuit)$ identifies with the exact sequence
\[
0 \To \lim_iM_i \To \prod_iM_i \To[\delta] \prod_iM_i \To {\lim_i}^1M_i \To 0
\]
of topological abelian groups, hence $\L(\lim_i^1\underline{M_i}) = \lim_i^1M_i$.
\end{remark}

\begin{lemma} \label{pro-vanishing-and-Mittag-Leffler--lemma}
For a tower $(M_i)_i$ of abelian groups the following are equivalent:
\begin{enumerate}
\item The tower is Mittag-Leffler and $\lim_iM_i=0$.
\item The tower is pro-zero, i.e.\ $\prolim_iM_i=0$ in $\Pro(\Ab)$.
\item It holds $\lim_i \underline{M_i} =0= \lim^1_i \underline{M_i}$ in $\Cond_\kappa(\Ab)$. 
\item It holds $\lim_i \underline{M_i} =0= \lim^1_i \underline{M_i}$ in $\Cond^\omega(\Ab)$.\footnote{The equivalence of this condition has been communicated to the authors by Yicheng Zhou.}
\end{enumerate}
\end{lemma}
\begin{proof}
First we show (1)$\Rightarrow$(2). The tower $(M_i)_i$ being Mittag-Leffler means that for every $i\geq 0$ the tower $(\im(M_j\to M_i))_{j\geq i}$ stabilises. Let $B_i\subset M_i$ be this stable image, i.e.\ $B_i=\bigcap_{j\geq i}\im(M_j\to M_i)$. 
Then $(B_i)_i$ is a subtower of $(M_i)_i$, hence $\lim_iB_i \subset \lim_iM_i =0$. Thus for every $i$ there exists a $j\geq i$ such that the map $M_j\to M_i$ is zero, hence the pro-system $\prolim_iM_i$ is zero. The direction (2)$\Rightarrow$(1) works analogously.

Now we assume (3). For any profinite set $S$ we have that $C(S,M_i)
= C(S,\ZZ) \otimes_\ZZ M_i$ and $C(S,\ZZ)$ is some large free abelian group \cite[Thm. 5.4]{condensed}, and so has $\bigoplus_\NN$ $\ZZ$ as a direct factor if $S$ is large enough. 
Thus there exists an extremally disconnected set $S$ such that 
\[
{\lim_i}^1 \bigoplus_\NN M_i \subset {\lim_i}^1C(S,M_i) = {\lim_i}^1 \underline{M_i}(S) \overset{(\clubsuit)}{=} \bigl( {\lim_i}^1\underline{M_i} \bigr)(S)=0
\]
where $(\clubsuit)$ holds as $\id_S$ is cofinal among all covers of $S$ (Remark~\ref{extremally-disconnected--rem}). The condition that $\lim_i^1\bigoplus_\NN M_i=0$ in $\Ab$ is equivalent to the tower $(M_i)_i$ being Mittag-Leffler \cite[Cor. 6]{emmanouil}, hence we have (1).

Assuming (1), for any extremally disconnected set $S$ we have 
\[
({\lim_i}^1\underline{M_i})(S) = {\lim_i}^1\underline{M_i}(S) = {\lim_i}^1C(S,M_i) \cong {\lim_i}^1\bigoplus_J M_i
\]
for some set $J$ and the last term vanishes as the tower is Mittag-Leffler. Hence $\lim_i^1\underline{M_i}$ in $\Cond_\kappa(\Ab)$.

The implication (3)$\Rightarrow$(4) follows since $j^*$ preserves limits and colimits.

It remains to show (4)$\Rightarrow$(1). For this let $S:=\NN\cup\{\infty\}\in\ProFin^\omega$. Then $\ZZ[S]$ is a projective object in $\Cond^\omega(\Ab)$\footnote{See Lecture 3 by Scholze in the lectures on Analytic Stacks, \url{https://youtu.be/me1KNo3WJHE?si=iSCb0CAYyCRvIn1R&t=2595}.} so that the equality $(\clubsuit)$ above holds for this particular $S$ in $\Cond^\omega(\Ab)$.
\end{proof}

\begin{example}
If a tower $(M_i)_i$ of abelian groups satisfies $\lim_iM_i=0=\lim_i^1M_i$, then it needs not to be pro-zero. For instance, consider the tower $(p^i\ZZ_p)_i$ for some prime number~$p$.
\end{example}




%%%%%%%%%%%%%%%%%%%%%%%%
\subsection{t-structures}
%%%%%%%%%%%%%%%%%%%%%%%%
In this subsection, we examine a t-structure on condensed spectra which is induced by the canonical t-structure on spectra.

\begin{lemma}
There exists a t-structure on $\Cond_\kappa(\Sp)$ with connective part $\Cond(\Sp)_{\geq 0} = \Cond(\Sp_{\geq 0})$ and coconnective part $\Cond(\Sp)_{\leq 0} = \Cond(\Sp_{\leq 0})$
\end{lemma}
\begin{proof}
%\hlcan{I think this follows follows from the proof of} \cite[1.3.2.7]{SAG} \hlcan{for light condensed spectra} 
Since $\Sp_{\geq 0}$ is presentable, this inherits to the presheaf category $\Fun(\EDS_{<\kappa}\op,\Sp_{\geq 0})$ \cite[5.5.3.6]{HTT} and its Bousfield localisation $\Cond_\kappa(\Sp_{\geq 0})$ \cite[5.5.4.15]{HTT}. The category $\Sp_{\geq 0}$ is closed under small colimits in $\Sp$ since the inclusion functor is left-adjoint. Since colimits in presheaf categories can be computed pointwise, $\Fun(\EDS_{<\kappa}\op,\Sp_{\geq 0})$ is closed under small colimits in $\Fun(\EDS_{<\kappa}\op,\Sp_{\geq 0})$. Since a colimit of sheaves is given by the sheafification of the colimits of the underlying presheaves, $\Cond_\kappa(\Sp_{\geq 0})$ is closed under small colimits in $\Cond_\kappa(\Sp)$. From the long exact sequences of homotopy groups one sees that $\Cond_\kappa(\Sp_{\geq 0})$ is closed under extension in $\Cond_\kappa(\Sp)$. These three closure properties imply the existence of a t-structure on $\Cond_\kappa(\Sp)$ whose connective part is $\Cond_\kappa(\Sp_{\geq 0})$ \cite[1.4.4.11]{HA}.

The category $\Cond_\kappa(\Sp)_{\geq 0}$ should be the same as the full subcategory of 0-connective objects in the sense of \cite[1.3.2.5]{SAG}. According to \cite[1.3.2.7]{SAG}, the coconnective part of the t-structure is given by the full subcategory of 0-truncated objects; these are precisely those sheaves that take value in 0-truncated spectra \cite[1.3.2.6]{SAG} (since the inclusion functor $\Sp_{\leq 0} \inj \Sp$ preserves limits. 
\end{proof}

\begin{lemma}
\label{lem-t-structure-pro}
There exists a t-structure on $\Pro(\Sp)$ with $\Pro(\Sp)_{\geq 0}$ given by the essential image of the functor $\Pro(\Sp_{\geq 0})\to\Pro(\Sp)$ and with $\Pro(\Sp)_{\leq 0}$ given by the full subcategory $\Pro(\Sp_{\leq 0}) \inj \Pro(\Sp)$.  

For $X=\prolim X_i\in\Pro(\Sp)$ and $n\in \ZZ$ we can compute the $n$-truncation pointwise, i.e. $\pi_n X\simeq \prolim \pi_n  X_i$.
\end{lemma}
\begin{proof}
First let us remark that $\Pro(\Sp)$ is a stable category \cite[Lem. 2.5]{kst-i}. The inclusion $\iota\colon\Sp_{\leq 0}\subseteq \Sp$ induces an adjunction 
$$
 \begin{tikzcd}
	\iota^{*}\colon \Pro(\Sp)\arrow[r,"",shift left = 0.3em]&\arrow[l,"",shift left = 0.3em]\Pro(\Sp_{\leq 0})\colon  \Pro(\iota)
\end{tikzcd}
$$
\cite[Ex. A.8.1.8]{SAG}. The functor $\Pro(\iota)$ is fully faithful \cite[Prop. A.8.1.9]{SAG}. 

We can apply the same argument to the truncation $\tau_{\geq 0}\colon \Sp\rightarrow \Sp_{\geq 0}$ and get an adjunction
$$
 \begin{tikzcd}
	\tau_{\geq 0}^{*}\colon \Pro(\Sp_{\geq 0})\arrow[r,"",shift left = 0.3em]&\arrow[l,"",shift left = 0.3em]\Pro(\Sp)\colon \Pro(\tau_{\geq 0}).
\end{tikzcd}
$$
Let us now define $\Pro(\Sp)_{\geq 0}$ as the essential image of $\tau_{\geq 0}^{*}$. We claim that the tuple 
$$
    (\Pro(\Sp)_{\geq 0}, \Pro(\Sp_{\leq 0}))
$$
defines a $t$-structure on $\Pro(\Sp)$.

Indeed, finite limits and colimit can be computed levelwise in $\Pro(\Sp)$ \cite[lem. 2.1]{kst-i}. So, checking the axioms of a $t$-structure reduces to the axioms for $(\Sp_{\geq 0},\Sp_{\leq 0})$, which holds by design.

For the homotopy groups note that $\tau_{\leq n}$ can be computed pointwise by construction. For $\tau_{\geq n}$ we note that for any $X\in \Pro(\Sp)$, we have a fibre sequence
$$
    \tau_{\geq n} X\to X\to \tau_{\leq n-1} X
$$
by the proof of \cite[Prop. 1.2.1.5]{HA}. As finite limits in $\Pro(\Sp)$ can be computed pointwise, we see that also $\tau_{\geq n}$ can be computed pointwise.
\end{proof}


\begin{prop}
The canonical functor $\gamma\colon \Pro(\Sp)\to\Cond(\Sp)$ is left t-exact.
\end{prop}
\begin{proof}
Since the inclusion functor $\Cond(\Sp_{\leq 0}) \inj \Cond(\Sp)$ preserves limits, the composition\linebreak $\Pro(\Sp_{\leq 0}) \to \Pro(\Sp) \to[\gamma] \Cond(\Sp)$ is equivalent to the composition $\Pro(\Sp_{\leq 0}) \to \Cond(\Sp_{\leq 0}) \inj \Cond(\Sp)$. Thus $\gamma$ is left t-exact. 
\end{proof}



%%%%%%%%%%%%%%%%%
\subsection{(Pre-)sheaves of condensed spaces and spectra}
%%%%%%%%%%%%%%%%%

In this subsection, we examine the relation between sheafification and stabilisation of presheaf categories. This will apply to the passage to condensed objects. First, we start with a general lemma:

\begin{lemma} \label{condensed-sheaves--lem}
Let $(\Ccal,\tau)$ and $(\Dcal,\sigma)$ be two sites. Then:
\begin{enumerate}
\item $\Fun(\Ccal\op,\Sh_\sigma(\Dcal)) \simeq \Sh_\sigma(\Dcal,\PSh(\Ccal))$
\item $\Sh_\tau(\Ccal,\Sh_\sigma(\Dcal)) \simeq \Sh_{\tau\times\sigma}(\Ccal\times\Dcal)$
\end{enumerate}
\end{lemma}
\begin{proof}
(i) Both categories are canonically subcategories of $\Fun(\Ccal\op\times\Dcal\op,\Spc)$. Since equivalences of presheaves can be computed pointwise, the condition for lying in both sides amounts to the sheaf condition for $\sigma$, hence the subcategories agree.\\
(ii) This holds since one can check the respective sheaf conditions in each factor independently.
\end{proof}

\begin{cor} \label{condensed-sheaves--cor}
Given a site $(\Ccal,\tau)$, there is a canonical equivalence
\[
\Sh_\tau(\Ccal,\Cond^{\omega}(\Spc)) \simeq \Cond^{\omega}(\Sh_\tau(\Ccal))
\]
between the categories of $\tau$-sheaves of condensed spaces and condensed $\tau$-sheaves.
The analogous result for $\Cond_\kappa(-)$ is true as well.
\end{cor}

\begin{lemma} \label{sheaves-and-stabilisation--lem}
Let $(\Ccal,\tau)$ be a site and $\Dcal$ be a category. There is a canonical equivalence
\[
\Sh_\tau(\Ccal,\Sp(\Dcal)) \simeq \Sp\bigl( \Sh(\Ccal,\Dcal) \bigr)
\]
between $\tau$-sheaves of spectra and spectrum objects in $\tau$-sheaves.
\end{lemma}
\begin{proof}
This is \cite[1.3.2.2]{SAG}.
\end{proof}

\begin{cor} \label{condensed-and-stabilisation--cor}
For any \infcat{} $\Dcal$ there is a canonical equivalence
\[
\Cond^{\omega}(\Sp(\Dcal)) \simeq \Sp(\Cond^{\omega}(\Dcal)).
\]
The analogous result for $\Cond_\kappa(-)$ is true as well.
\end{cor}
\begin{proof}
This follows from Lemma~\ref{sheaves-and-stabilisation--lem}.
\end{proof}





%%%%%%%%%%%%%%%%%%%%%%%%%%%%%%%%%%%%%%%%%%%%%%%%%
\subsection*{Internal mapping objects}
%%%%%%%%%%%%%%%%%%%%%%%%%%%%%%%%%%%%%%%%%%%%%%%%%

In order to render the functor $\Kan_{\geq 0} \colon \Rig_k \to \Pro(\Spc) \to \Cond^{\omega}(\Spc)$ representable, we need to have internal condensed mapping spectra in $\RigH(k,\Cond^{\omega}(\Spc))$. ByRemark \ref{rem-RigH-monoidal} it is enough to see that the $\otimes$-product in $\Cond^{\omega}(\Spc)$ given by Lemma \ref{condensed-objects-symmetric-monoidal} preserves colimits in each variable. 

\begin{lem}
    The cartesian symmetric monoidal structure on $\Cond(\Spc)$ is closed, i.e.\ for any objects $Y,Z\in\Cond(\Spc)$ there exists an internal mapping object $\Mapu(Y,Z)\in\Cond(\Spc)$ such that for any $X\in\Cond(\Spc)$ there is a canonical equivalence 
\[
\Map_{\Cond(\Spc)}(X,\Mapu(Y,Z)) \simeq \Map_{\Cond(\Spc)}(X\times Y,Z).
\]
Setting $X=y(S)$ for $S\in\EDS$ we get that
\[
\Mapu(Y,Z)(S) \simeq \Map_{\Cond(\Spc)}(y(S),\Mapu(Y,Z)) \simeq \Map_{\Cond(\Spc)}(y(S)\times Y,Z).
\]

Moreover, for any $\Vcal\in \CAlgPr$, we have $\Cond^{\omega}(\Vcal)$ admits a closed monoidal structure.
\end{lem}
\begin{proof}
    This follows immediately from the definition, since $\Cond^{\omega}(\Spc)$ as an $\infty$-topos and thus colimits are universal. If $\Vcal$ is a symmetric monoidal presentable category, then we have $\Cond^{\omega}(\Spc)\otimes \Vcal\simeq \Cond^{\omega}(\Vcal)$ \cite[Prop. 2.4]{DrewMHM}.
\end{proof}



%%%%%%%%%%%%%%%%%%%%%%%%%%%%%%%%%%%%%%%%%%%%%
\section{Enriched Yoneda}
\label{app.yoneda}
%%%%%%%%%%%%%%%%%%%%%%%%%%%%%%%%%%%%%%%%%%%%%

In this section, we want to recall the enriched Yoneda lemma in the case that is of interest for us, following Hinich \cite{Hinich}. We will give a short summary of the constructions.

In the following, we will fix a category $\Ccal$ and a symmetric monoidal presentable category $\Ecal$. Note that the work of Hinich works in greater generality but we restrict to this case.

\begin{notation}
    We denote by $Y\colon \Ccal\op\times \Ccal\rightarrow \Spc$ the functor $(x,y)\mapsto \Hom(x,y)$. We denote the composition of this functor with the unit $\Spc\rightarrow \Ecal$ by $\Ytilde$.
\end{notation}

In \cite[\S 3.1.3]{Hinich} Hinich defines a category of quivers $\textup{Quiv}_{\Ccal}^{\textup{BM}}(\Ecal)$ as an $\infty$-operad\footnote{In \textit{loc.\! cit.} Hinich uses a different notion of operads, based on strong approximations of $\infty$-operads. We will ignore this detail as this section stays valid, having \cite[\S 2.7]{Hinich} in mind.} fibred over the $\infty$-operad of bimodules $\textup{BM}$ (the $\infty$-operad with colours $\lbrace a_{+},m,a_{-}\rbrace$). The colours of $\textup{Quiv}_{\Ccal}^{\textup{BM}}(\Ecal)$ can be described as
\begin{enumerate}
    \item[$\bullet$] $\textup{Quiv}_{\Ccal}(\Ecal)\coloneqq\textup{Quiv}_{\Ccal}(\Ecal)_{a_{+}} \simeq \Fun(\Ccal^{op}\times \Ccal,\Ecal)$,
    \item[$\bullet$]  $\textup{Quiv}_{\Ccal}(\Ecal)_{m} \simeq \Fun(\Ccal,\Ecal)$, and 
    \item[$\bullet$] $\textup{Quiv}_{\Ccal}(\Ecal)_{a_{-}}\simeq\Ecal$.
\end{enumerate}

Let us remark that since $\Ecal$ is symmetric monoidal, we have an equivalence categories 
$$
    \textup{Quiv}_{\Ccal}(\Ecal)\simeq \textup{Quiv}_{\Ccal\op}(\Ecal).
$$
In particular, the functor $\Ytilde$ defines an associative algebra object $\Ytilde\op$ inside $\textup{Quiv}_{\Ccal\op}(\Ecal)$ via the above equivalence in the sense of \cite{Hinich}. Thus, $\Ytilde\op$ acts on $\Fun(\Ccal\op,\Ecal)$ from the left. 
\begin{defi}[\protect{\cite[\S 6.2.2]{Hinich}}]
    We define the category of $\Ecal$-presheaves on $\Ccal$ as 
    $$
        P_{\Ecal}(\Ytilde) \coloneqq \textup{LMod}_{\Ytilde\op}(\Fun(\Ccal\op,\Ecal)).
    $$
\end{defi}
There is a lot to digest in this definition of $\Ecal$-presheaves. But in our context we can greatly simplify the definition. 

First note $\textup{Quiv}_{\Ccal\op}(\Ecal)\simeq \textup{Quiv}_{\Ccal}(\Ecal)^{\textup{rev}}$,
where the superscript $\textup{rev}$ denotes the reversed monoidal category in the sense of \cite[Rem. 4.1.1.7]{HA} (this follows from the construction \cite[\S 6.2.1]{Hinich}). Further, by \cite[\S 4.7.3]{Hinich}, we have that $\Ytilde$ is the unit object inside  $\textup{Quiv}_{\Ccal}(\Ecal)$ and thus $Y\op$ is the unit object in $\textup{Quiv}_{\Ccal}(\Ecal)^{\textup{rev}}$. In particular, we obtain
\begin{equation}\tag{$\clubsuit$}
    \textup{LMod}_{Y\op}(\Fun(\Ccal\op,\Ecal))\simeq \Fun(\Ccal\op,\Ecal)
\end{equation}
under the above identifications \cite[Prop. 4.2.4.9]{HA}.
By design of $\textup{Quiv}_{\Ccal\op}^{\textup{BM}}(\Ecal)$, we get an induced right $\Ecal^{\textup{rev}}$ action on $\Fun(\Ccal\op,\Ecal)$, which is the same as a left $\Ecal$-action. This action is induced by the pointwise monoidal structure.

The left $\Ecal$-module structure on $\Fun(\Ccal\op,\Ecal)$ induces an enrichment as the $\otimes$-product of $\Ecal$ commutes with arbitrary colimits, i.e. we obtain an internal mapping object\linebreak $\Homline_{\Ecal}(F,G)\in \Ecal$ for all $F,G\in \Fun(\Ccal\op,\Ecal)$. By \cite[\S 6.2.4]{Hinich}, we obtain a map $Y'\colon \Ccal\rightarrow \textup{LMod}_{Y\op}(\Fun(\Ccal\op,\Ecal))$ with a compatible system of maps 
$$
    \Ytilde(x,y) \otimes Y'(x)\rightarrow Y'(y)
$$
for all $x,y\in \Ccal$. Under the equivalence ($\clubsuit$), we have $Y'\simeq \Ytilde$ \cite[\S 4.7.4]{Hinich} (see also the proof of \cite[6.2.6 Cor.]{Hinich}). This allows us to formulate an enriched Yoneda lemma.

\begin{lem}[Enriched Yoneda \protect{\cite[\S 6.2.7]{Hinich}}]
\label{lem-yoneda}
    Let $F\in \Fun(\Ccal\op,\Ecal)$ be an $\Ecal$-presheaf on $\Ccal$. Then, we have an equivalence 
    $$
        F(x)\simeq \Homline_{\Ecal}(\Ytilde(x),F)
    $$
    for all $x\in \Ccal$. Moreover, the induced functor $\Ytilde\colon \Ccal\rightarrow \Fun(\Ccal\op,\Ecal)$ is fully faithful. 
\end{lem}

\end{appendix}





%------------------------------------------------------------------

%------------------------------------------------------------------

%==================================================================

\bibliographystyle{alphaurl}
\bibliography{RigidSH}


%--------------------------------
\bigskip
  \footnotesize

  Christian Dahlhausen:\\  \textsc{Institut für Mathematik, Universtität Heidelberg, Im Neuenheimer Feld 205, 69120 Heidelberg, Germany}\par\nopagebreak
  \textit{E-mail address:}  \texttt{cdahlhausen@mathi.uni-heidelberg.de}

  \medskip

  Can Yaylali:\\ \textsc{Fachbereich Mathematik, Technische Universität Darmstadt, Schlossgartenstr. 7, 64289 Darmstadt, Germany}\par\nopagebreak
  \textit{E-mail address:} \texttt{yaylali@mathematik.tu-darmstadt.de}
%--------------------------------


\end{document}

\newcommand{\etalchar}[1]{$^{#1}$}
\begin{thebibliography}{DL{\O}{\etalchar{+}}07}

\bibitem[Abb10]{egr}
Ahmed Abbes.
\newblock {\em {\'E}l\'ements de g\'eom\'etrie rigide. {V}olume {I}}, volume
  286 of {\em Progress in Mathematics}.
\newblock Birkh\"auser/Springer Basel AG, Basel, 2010.
\newblock Construction et \'etude g\'eom\'etrique des espaces rigides.

\bibitem[AGV22]{AGV}
Joseph Ayoub, Martin Gallauer, and Alberto Vezzani.
\newblock The six-functor formalism for rigid analytic motives.
\newblock {\em Forum Math. Sigma}, 10:Paper No. e61, 182, 2022.
\newblock \href {https://doi.org/10.1017/fms.2022.55}
  {\path{doi:10.1017/fms.2022.55}}.

\bibitem[And23]{andreychev-thesis}
Grigory Andreychev.
\newblock {\em K-Theorie adischer Räume}.
\newblock PhD thesis, Rheinische Friedrich-Wilhelms-Universität Bonn,
  September 2023.
\newblock \url{https://nbn-resolving.org/urn:nbn:de:hbz:5-72174}.
\newblock URL: \url{https://hdl.handle.net/20.500.11811/11040}.

\bibitem[Ayo07]{AyoubThesis}
Joseph Ayoub.
\newblock Les six op\'{e}rations de {G}rothendieck et le formalisme des cycles
  \'{e}vanescents dans le monde motivique. {(I, II)}.
\newblock {\em Ast\'{e}risque}, (314, 315), 2007.

\bibitem[Ayo15]{AyoubRig}
Joseph Ayoub.
\newblock Motifs des vari\'{e}t\'{e}s analytiques rigides.
\newblock {\em M\'{e}m. Soc. Math. Fr. (N.S.)}, (140-141):vi+386, 2015.
\newblock \href {https://doi.org/10.24033/msmf.449}
  {\path{doi:10.24033/msmf.449}}.

\bibitem[CD19]{CD1}
Denis-Charles Cisinski and Fr\'{e}d\'{e}ric D\'{e}glise.
\newblock {\em Triangulated categories of mixed motives}.
\newblock Springer Monographs in Mathematics. Springer, Cham, [2019] \copyright
  2019.
\newblock \href {https://doi.org/10.1007/978-3-030-33242-6}
  {\path{doi:10.1007/978-3-030-33242-6}}.

\bibitem[Cis13]{cisinski-descente}
Denis-Charles Cisinski.
\newblock Descente par \'eclatements en {$K$}-th\'eorie invariante par
  homotopie.
\newblock {\em Ann. of Math. (2)}, 177(2):425--448, 2013.
\newblock \href {https://doi.org/10.4007/annals.2013.177.2.2}
  {\path{doi:10.4007/annals.2013.177.2.2}}.

\bibitem[Con07]{NagataComp}
Brian Conrad.
\newblock Deligne's notes on {N}agata compactifications.
\newblock {\em J. Ramanujan Math. Soc.}, 22(3):205--257, 2007.

\bibitem[Dah24]{dahli-thesis-paper}
Christian Dahlhausen.
\newblock Continuous {K}-theory and cohomology of rigid spaces.
\newblock {\em Manuscripta Math.}, 173(1-2):119--153, 2024.
\newblock \href {https://doi.org/10.1007/s00229-023-01470-x}
  {\path{doi:10.1007/s00229-023-01470-x}}.

\bibitem[DL{\O}{\etalchar{+}}07]{MotCoh}
Bj{\o}rn Dundas, Marc Levine, Paul {\O}stv{\ae}r, Oliver R\"ondigs, and
  Vladimir Voevodsky.
\newblock {\em Motivic Homotopy Theory: Lectures at a Summer School in
  Nordfjordeid, Norway, August 2002}.
\newblock 01 2007.
\newblock \href {https://doi.org/10.1007/978-3-540-45897-5}
  {\path{doi:10.1007/978-3-540-45897-5}}.

\bibitem[Dre18]{DrewMHM}
Brad Drew.
\newblock Motivic hodge modules, 2018.
\newblock \href {http://arxiv.org/abs/1801.10129} {\path{arXiv:1801.10129}}.

\bibitem[Emm96]{emmanouil}
Ioannis Emmanouil.
\newblock Mittag-{L}effler condition and the vanishing of {$\varprojlim^1$}.
\newblock {\em Topology}, 35(1):267--271, 1996.
\newblock \href {https://doi.org/10.1016/0040-9383(94)00056-5}
  {\path{doi:10.1016/0040-9383(94)00056-5}}.

\bibitem[FK18]{fuji-kato}
Kazuhiro Fujiwara and Fumiharo Kato.
\newblock {\em Foundations of {R}igid {G}eometry {I}}, volume~7 of {\em
  Monographs in Mathematics}.
\newblock EMS, 2018.

\bibitem[GW10]{GW1}
Ulrich G\"ortz and Torsten Wedhorn.
\newblock {\em Algebraic {G}eometry {I}}.
\newblock Advanced Lectures in Mathematics. Vieweg + Teubner, Wiesbaden, 2010.
\newblock Schemes with examples and exercises.
\newblock URL: \url{https://doi.org/10.1007/978-3-8348-9722-0}.

\bibitem[Hin20]{Hinich}
Vladimir Hinich.
\newblock Yoneda lemma for enriched $\infty$-categories.
\newblock {\em Advances in Mathematics}, 367:107129, 2020.
\newblock URL:
  \url{https://www.sciencedirect.com/science/article/pii/S0001870820301559},
  \href {https://doi.org/https://doi.org/10.1016/j.aim.2020.107129}
  {\path{doi:https://doi.org/10.1016/j.aim.2020.107129}}.

\bibitem[Hov01]{Hovey}
Mark Hovey.
\newblock Spectra and symmetric spectra in general model categories.
\newblock {\em J. Pure Appl. Algebra}, 165(1):63--127, 2001.
\newblock \href {https://doi.org/10.1016/S0022-4049(00)00172-9}
  {\path{doi:10.1016/S0022-4049(00)00172-9}}.

\bibitem[Hoy17]{Hoy-EQ}
Marc Hoyois.
\newblock The six operations in equivariant motivic homotopy theory.
\newblock {\em Adv. Math.}, 305:197--279, 2017.
\newblock \href {https://doi.org/10.1016/j.aim.2016.09.031}
  {\path{doi:10.1016/j.aim.2016.09.031}}.

\bibitem[Hoy24]{Hoy-EG-arxiv-v5}
Marc Hoyois.
\newblock The six operations in equivariant motivic homotopy theory.
\newblock \href{https://arxiv.org/abs/1509.02145v5}{arXiv:1509.02145v5}, 2024.

\bibitem[Hü24]{huebner-adic}
Katharina Hübner.
\newblock Adic spaces.
\newblock Lectures at the spring school `Nonarchimedean Geometry and
  Eigenvarieties' in March 2023 in Heidelberg.
  \href{https://arxiv.org/abs/2405.06435}{arXiv:2405.06435}, 2024.

\bibitem[Isa02]{isaksen-calculating}
Daniel~C. Isaksen.
\newblock Calculating limits and colimits in pro-categories.
\newblock {\em Fund. Math.}, 175(2):175--194, 2002.
\newblock \href {https://doi.org/10.4064/fm175-2-7}
  {\path{doi:10.4064/fm175-2-7}}.

\bibitem[KST19a]{kst-i}
Moritz Kerz, Shuji Saito, and Georg Tamme.
\newblock K-theory of non-archimedean rings. {I}.
\newblock {\em Doc. Math.}, 24:1365–1411, 2019.
\newblock \href {https://doi.org/10.25537/dm.2019v24.1365-1411}
  {\path{doi:10.25537/dm.2019v24.1365-1411}}.

\bibitem[KST19b]{kst-bass-quillen}
Moritz Kerz, Shuji Saito, and Georg Tamme.
\newblock Towards a non-archimedean analytic analog of the {B}ass-{Q}uillen
  conjecture.
\newblock {\em Journal of the Institute of Mathematics of Jussieu}, pages
  1--16, 2019.
\newblock \href {https://doi.org/10.1017/S147474801900001X}
  {\path{doi:10.1017/S147474801900001X}}.

\bibitem[KST23]{kst-ii}
Moritz Kerz, Shuji Saito, and Georg Tamme.
\newblock {$K$}-theory of non-{A}rchimedean rings {II}.
\newblock {\em Nagoya Math. J.}, 251:669--685, 2023.
\newblock \href {https://doi.org/10.1017/nmj.2023.4}
  {\path{doi:10.1017/nmj.2023.4}}.

\bibitem[Lur09]{HTT}
Jacob Lurie.
\newblock {\em Higher topos theory}, volume 170 of {\em Annals of Mathematics
  Studies}.
\newblock Princeton University Press, Princeton, NJ, 2009.
\newblock \href {https://doi.org/10.1515/9781400830558}
  {\path{doi:10.1515/9781400830558}}.

\bibitem[Lur17]{HA}
Jacob Lurie.
\newblock Higher algebra.
\newblock \url{https://www.math.ias.edu/~lurie/papers/HA.pdf}, 2017.

\bibitem[Lur18]{SAG}
Jacob Lurie.
\newblock Spectral algebraic geometry.
\newblock \url{https://www.math.ias.edu/~lurie/papers/SAG-rootfile.pdf}, 2018.

\bibitem[LZ12]{Liu-Zheng}
Yifeng Liu and Weizhe Zheng.
\newblock Enhanced six operations and base change theorem for higher artin
  stacks, 2012.
\newblock arXiv:1211.5948.
\newblock URL: \url{https://arxiv.org/abs/1211.5948}, \href
  {https://doi.org/10.48550/ARXIV.1211.5948}
  {\path{doi:10.48550/ARXIV.1211.5948}}.

\bibitem[Man22]{Mann-Thesis}
Lucas Mann.
\newblock A $p$-adic 6-functor formalism in rigid-analytic geometry, 2022.
\newblock \href {http://arxiv.org/abs/2206.02022} {\path{arXiv:2206.02022}}.

\bibitem[Mor16]{Morrow}
Matthew Morrow.
\newblock A historical overview of pro cdh descent in algebraic $k$-theory and
  its relation to rigid analytic varieties, 2016.
\newblock \href {http://arxiv.org/abs/1612.00418} {\path{arXiv:1612.00418}}.

\bibitem[MV99]{MV1}
Fabien Morel and Vladimir Voevodsky.
\newblock {${\bf A}^1$}-homotopy theory of schemes.
\newblock {\em Inst. Hautes \'{E}tudes Sci. Publ. Math.}, (90):45--143 (2001),
  1999.
\newblock URL: \url{http://www.numdam.org/item?id=PMIHES_1999__90__45_0}.

\bibitem[PPR09]{PPR-KGL}
Ivan Panin, Konstantin Pimenov, and Oliver R\"ondigs.
\newblock On {V}oevodsky's algebraic {$K$}-theory spectrum.
\newblock In {\em Algebraic topology}, volume~4 of {\em Abel Symp.}, pages
  279--330. Springer, Berlin, 2009.
\newblock URL: \url{https://doi.org/10.1007/978-3-642-01200-6_10}, \href
  {https://doi.org/10.1007/978-3-642-01200-6\_10}
  {\path{doi:10.1007/978-3-642-01200-6\_10}}.

\bibitem[Ray74]{raynaud}
Michel Raynaud.
\newblock G\'eom\'etrie analytique rigide d'apr\`es {T}ate, {K}iehl,{$\cdots
  $}.
\newblock In {\em Table {R}onde d'{A}nalyse {N}on {A}rchim\'edienne ({P}aris,
  1972)}, volume Tome 102 of {\em Suppl\'ement au Bull. Soc. Math. France},
  pages 319--327. Soc. Math. France, Paris, 1974.
\newblock \href {https://doi.org/10.24033/msmf.170}
  {\path{doi:10.24033/msmf.170}}.

\bibitem[Rob14]{RobaloDiss}
Marco Robalo.
\newblock Th\'eorie homotopique motivique des espaces non-commutatifs.
\newblock 2014.
\newblock PhD-thesis, University of Montpellier.

\bibitem[Rob15]{Robalo-Thesis}
Marco Robalo.
\newblock {$K$}-theory and the bridge from motives to noncommutative motives.
\newblock {\em Adv. Math.}, 269:399--550, 2015.
\newblock \href {https://doi.org/10.1016/j.aim.2014.10.011}
  {\path{doi:10.1016/j.aim.2014.10.011}}.

\bibitem[Sch19]{condensed}
Peter Scholze.
\newblock Lectures on {C}ondensed {M}athematics.
\newblock available online:
  \href{https://www.math.uni-bonn.de/people/scholze/Condensed.pdf}{https://www.math.uni-bonn.de/people/scholze/Condensed.pdf},
  2019.

\bibitem[Sig16]{sigloch-thesis}
Helene Sigloch.
\newblock {\em Homotopy Theory for Rigid Analytic Varieties}.
\newblock PhD thesis, Albert-Ludwigs-Universit¨at Freiburg im Breisgau, March
  2016.
\newblock \url{https://freidok.uni-freiburg.de/data/11742}.
\newblock URL: \url{https://freidok.uni-freiburg.de/data/11742}.

\bibitem[{Sta}23]{stacks-project}
The {Stacks Project Authors}.
\newblock \textit{Stacks Project}.
\newblock \url{https://stacks.math.columbia.edu}, 2023.

\bibitem[SW20]{weinstein-scholze}
Peter Scholze and Jared Weinstein.
\newblock {\em Berkeley lectures on {$p$}-adic geometry}, volume 207 of {\em
  Annals of Mathematics Studies}.
\newblock Princeton University Press, Princeton, NJ, 2020.

\bibitem[TT90]{tt90}
Robert~W. Thomason and Thomas Trobaugh.
\newblock Higher algebraic {$K$}-theory of schemes and of derived categories.
\newblock In {\em The {G}rothendieck {F}estschrift, {V}ol.\ {III}}, volume~88
  of {\em Progr. Math.}, pages 247--435. Birkh\"auser Boston, Boston, MA, 1990.
\newblock URL: \url{http://dx.doi.org/10.1007/978-0-8176-4576-2_10}, \href
  {https://doi.org/10.1007/978-0-8176-4576-2_10}
  {\path{doi:10.1007/978-0-8176-4576-2_10}}.

\bibitem[Voe98]{voevodsky-icm-1998}
Vladimir Voevodsky.
\newblock {${\bf A}^1$}-homotopy theory.
\newblock In {\em Proceedings of the {I}nternational {C}ongress of
  {M}athematicians, {V}ol. {I} ({B}erlin, 1998)}, pages 579--604, 1998.

\bibitem[Wei13]{k-book}
Charles~A. Weibel.
\newblock {\em The {$K$}-book}, volume 145 of {\em Graduate Studies in
  Mathematics}.
\newblock American Mathematical Society, Providence, RI, 2013.
\newblock An Introduction to Algebraic $K$-theory.

\end{thebibliography}
